\documentclass[11pt]{article}
\usepackage[margin=1in]{geometry}
\usepackage{amsmath,amssymb,amsthm,mathtools,booktabs,microtype,array}
\usepackage[colorlinks=true,linkcolor=blue,citecolor=blue,urlcolor=blue]{hyperref}
\hypersetup{pdftitle={A zero-free strip for the L-functions of the elliptic curves y\textasciicircum 2=x\textasciicircum 3+d: the 11/12 method and its extension to 7/8},pdfauthor={OpenAI Astra-6 and Claude Opus 5.5, at the behest of Mateo Restrepo Borrero and Michael O. Rubinstein},pdfsubject={L(E\textunderscore d,s) is nonzero for Re s > 17/12 (Part I) and Re s > 11/8 (Part II), where E\textunderscore d: y\textasciicircum 2=x\textasciicircum 3+d}}
\allowdisplaybreaks[2]
\setlength{\emergencystretch}{2em}
\numberwithin{equation}{section}
\newcommand{\OO}{\mathbb Z[\omega]}
\newcommand{\N}{\mathrm N}
\newcommand{\muk}{\mu_K}
\newcommand{\ind}{\mathbf1}
\newcommand{\E}{\mathcal E}
\newcommand{\ce}{\check e}
\newcommand{\ssum}{\sideset{}{^*}\sum}
\newcommand{\Tr}{\operatorname{Tr}}
\newtheorem{proposition}{Proposition}
\newtheorem{theorem}{Theorem}
\newtheorem{lemma}{Lemma}
\title{A zero-free strip for the $L$-functions\\of the elliptic curves $y^2=x^3+d$\\[1ex]
{\large the $11/12$ method and its extension to $7/8$}}
\author{OpenAI Astra-6 \quad and \quad Claude Opus 5.5\\[1ex]
{\normalsize at the behest of Mateo Restrepo Borrero and Michael O.\ Rubinstein}}
\date{}
\begin{document}
\maketitle
\begin{center}\setlength{\fboxsep}{6pt}
\fbox{\parbox{0.9\linewidth}{\small\textbf{Working notes.} These are working notes. They describe an argument
that is still being developed, and the details have yet to be checked carefully.}}
\end{center}

\begin{abstract}
For a nonzero integer $d$ let $E_d$ be the elliptic curve
$y^2=x^3+d$ over $\mathbb Q$. These are exactly the curves with
$j$-invariant $0$, and each has complex multiplication by $\OO$.
Part~I of these notes explains how OpenAI's $11/12$ argument adapts to
give the following statement: for every $d\ne0$,
\[
 L(E_d,s)\ne0\qquad\text{for }\Re s>17/12 .
\]
The critical strip of $L(E_d,s)$ is $1/2<\Re s<3/2$, so this is a
zero-free strip $17/12<\Re s<3/2$ along its right edge. In the analytic
normalization, with critical strip $0<\Re s<1$, it says that
$L(E_d,s+\frac12)\ne0$ for $\Re s>11/12$, the same boundary as OpenAI's
$11/12$ theorem for $\zeta(s)$. Since $L(E_d,s)$ is the $L$-function of a
Hecke character of $\mathbb Q(\sqrt{-3})$ of infinite order, the statement
is not covered by OpenAI's $7/8$ theorem, which concerns
finite-order Hecke $L$-functions over $\mathbb Q(\sqrt{-3})$. We begin with
ordinary integer coefficients obtained by counting points modulo primes,
write out the mean squares, and isolate the sixth-power rows that pick
out the chosen curve.

Part~II (Sections~\ref{sec:78}--\ref{sec:78status}) adapts OpenAI's $7/8$
argument to the same curves. It outlines a proof, conditional on restatements
listed in Section~\ref{sec:78status}, that $L(E_d,s)\ne0$ for $\Re s>11/8$,
that is, of a zero-free strip $11/8<\Re s<3/2$. More generally, it outlines a
proof that $L(s,\overline\alpha\eta)\ne0$ for $\Re s>7/8$ for every finite-order Hecke
character $\eta$ of $\mathbb Q(\sqrt{-3})$, where $\alpha(n)=n/|n|$. The
adaptation changes three things: the weight of the target, the
bookkeeping of the characters attached to rows, and an angular factor
in the moment estimates. We list the statements of OpenAI's paper that
must still be written out in this form. Section~\ref{sec:higherweight}
explains why the same method should extend to CM newforms of every weight
with complex multiplication by $\mathbb Q(\sqrt{-3})$, such as $\eta(3z)^8$ and
$\eta(2z)^3\eta(6z)^3$: weight $k+1$ needs $k+1$ derivatives of the cubic theta
function and, as far as we can see, nothing else new. This would also give the zero-free half-plane
for every symmetric power $L$-function of $E_d$.
\end{abstract}

\paragraph{The method in six sentences.}
\begin{enumerate}
\item \emph{The sum to bound.} Up to finitely many Euler factors,
 $1/L(E_d,s+\frac12)=\sum_mq_d(m)m^{-s}$, where $q_d$ is the multiplicative function with
 $q_d(p)=-a_p(d)/\sqrt p$, $q_d(p^2)=1$ and $q_d(p^j)=0$ for $j\ge3$ ($p\nmid6d$). So a bound
 $S_d(D)=\sum_mq_d(m)W(m/D)\ll D^{11/12+\varepsilon}$ for smooth $W$ gives $L(E_d,s)\ne0$ for
 $\Re s>17/12$, exactly as the same bound for the coefficients $c(m)$ of
 $1/(\zeta(s)L(s,\chi_{-3}))$ gives $\zeta(s)\ne0$ for $\Re s>11/12$.
\item \emph{The same sum over Eisenstein integers.} Grouping by norms,
 \[
  S_d(D)=\sum_n\muk(n)\,\frac n{|n|}\,\chi_n\bigl((4d)^5\bigr)W(\N(n)/D)
 \]
 over primary $n\in\OO$ coprime to $6$, while the zeta sum is $\sum_n\muk(n)W(\N(n)/D)$. The curve enters only through the unimodular
 factor $n/|n|$ and the sextic residue symbol of the fixed number $(4d)^5$; together these encode the
 point counts through $a_p(d)=\Tr\bigl(\pi\overline{\chi_\pi(4d)}\bigr)$.
\item \emph{What is averaged.} Replacing $(4d)^5$ by an arbitrary Eisenstein integer $u$ gives the
 family $A_u(D)=\sum_n\muk(n)\frac n{|n|}\chi_n(u)W(\N(n)/D)$, and the quantity bounded is the mean square
 $\sum_{0<\N(u)\le H}|A_u(D)|^2$ over all $u$ with $\N(u)\le H=D^{1+\vartheta}$. That is about $H$ rows
 of length about $D$, and the method shows the mean square is $\ll D^{2+\vartheta+\varepsilon}=D^{1+\varepsilon}H$, the
 size given by square-root cancellation in each row. The zeta argument averages exactly the same rows
 without the factor $n/|n|$.
\item \emph{Isolating the curve.} Since $\chi_n(u\pi^6)=\chi_n(u)$ whenever $\pi\nmid n$, the roughly
 $H^{1/6}/\log H$ rows $u=(4d)^5\pi^6$, with $\pi$ prime and $\N(\pi)\asymp H^{1/6}$, all equal $S_d(D)$ up to
 $O(D/\N(\pi))$. Positivity of the mean square then gives $H^{1/6}|S_d(D)|^2\lesssim DH$, i.e.\
 $S_d(D)\lesssim D^{11/12}$ when $H\approx D$. For zeta the rows $\pi^6$ play the same role for
 $\sum_n\muk(n)W(\N(n)/D)$.
\item \emph{Poisson summation and the theta function.} Expanding the square and applying Poisson
 summation in $u$ turns the paired Möbius coefficients into normalized cubic Gauss sums times
 $(\bar n/|n|)^2$ (for zeta, times $\bar n/|n|$). Sums of these come from Kubota's cubic theta function
 differentiated twice in $\bar z$ (for zeta, once), and its symmetry turns the sextic twists into
 quadratic residue symbols, which the quadratic large sieve controls.
\item \emph{Closing the argument.} The theta function also brings in terms at cube indices $nb^3$.
 Möbius inversion over cubes and two further Poisson summations reduce the number of rows by a fixed
 power of $D$ at each step. This induction, identical to the zeta case once the powers of $n/|n|$ are
 tracked, proves the mean-square bound, hence the exponent $11/12$ and the arithmetic boundary $17/12$.
\end{enumerate}

\paragraph{What is proved here, and what is quoted.}
These notes follow the organization of the zeta guide~\cite{ZG}, adapting
the method of~\cite{OAI1112}. The elementary deduction from the CM mean-square
estimate is given in full. We state the harder analytic estimates precisely
and explain their mechanism and how the arguments of~\cite{OAI1112} adapt;
we do not reproduce complete proofs of the theta transformation or the
descent estimates. Here ``CM curves'' means the $j=0$ curves over $\mathbb Q$.
Other CM fields would require their own arithmetic and analytic
inputs. Throughout, the curve parameter $d$ is fixed; constants may depend
on it. Part~II adapts~\cite{OAI78};
Section~\ref{sec:78status} says what is proved there and what is carried
over from~\cite{OAI78} by inspection.

\paragraph{Higher weight.} The angular factor $\alpha$ enters the theta step
only through the number of derivatives, and elsewhere only through Gamma
factors with the same growth. A CM newform of weight
$k+1$ with complex multiplication by $\mathbb Q(\sqrt{-3})$ corresponds to the
character $\overline\alpha^{\,k}\eta$ and needs $k+1$ derivatives. For each fixed $k$,
the formulas have the same structure and give the same length exponents; constants may depend on $k$.
Section~\ref{sec:higherweight}
explains why both the $11/12$ and the $7/8$ boundaries should therefore extend
to these forms, for example to $\eta(3z)^8$ (weight $4$) and $\eta(2z)^3\eta(6z)^3$
(weight $3$), and hence, by factoring, to the symmetric power $L$-functions of
$E_d$. This is an expectation, not a claim.

\section{Why start with \texorpdfstring{$11/12$}{11/12} rather than \texorpdfstring{$7/8$}{7/8}?}

For this family, the $11/12$ method is substantially simpler to present.
Its main analytic task is one positive mean square, and the passage from
that bound to the zero-free region takes only a few inequalities and a
Mellin integral. Relative to the zeta argument, the main new ingredient is
the angular factor $\alpha(n)=n/|n|$: one theta derivative becomes two.

The $7/8$ argument~\cite{OAI78,G78} has additional layers: a supremum of
zero locations over a whole family, an average evaluated in two ways,
short prime sums arranged to cancel an unwanted term, and two different
moment estimates to count exceptional rows. These are substantial additions
to the machinery used below.

There is also a scope issue. The stated $7/8$ theorem in~\cite{OAI78}
concerns \emph{finite-order} Hecke characters over $\mathbb Q(\sqrt{-3})$.
The characters for our curves include $n/|n|$ and are not of finite order.
Thus $7/8+1/2=11/8$ is the corresponding arithmetic boundary only
after the required CM adaptation; it is not a consequence of
shifting the stated finite-order theorem. Part~I (Sections~\ref{sec:integer}--\ref{sec:descent})
explains the CM $11/12$ argument. Part~II
(Sections~\ref{sec:78}--\ref{sec:78status}) carries out the $7/8$
adaptation.

\section{Start with ordinary integers and point counts}\label{sec:integer}

At a prime $p\nmid6d$, put
\[
 a_p(d)=p+1-\#E_d(\mathbb F_p)
       =-\sum_{x=0}^{p-1}\left(\frac{x^3+d}{p}\right),
 \qquad \lambda_d(p)=\frac{a_p(d)}{\sqrt p}.
\]
The symbol on the right is the ordinary Legendre symbol, with value zero
on multiples of $p$. Indeed, the number of solutions in $y$ to
$y^2=x^3+d$ is $1+((x^3+d)/p)$, and there is one point at infinity.
The good Euler factor is
\[
 L_p(E_d,s+\tfrac12)^{-1}
 =1-\lambda_d(p)p^{-s}+p^{-2s}.
\]
The shift is chosen so that $L(E_d,s+1/2)$ has critical line
$\Re s=1/2$; the arithmetic function $L(E_d,s)$ has critical line
$\Re s=1$.
Consequently the reciprocal coefficients away from $6d$ are the
multiplicative function $q_d$ specified by
\begin{equation}\label{eq:qd}
 q_d(1)=1,\qquad q_d(p)=-\lambda_d(p),\qquad
 q_d(p^2)=1,\qquad q_d(p^j)=0\quad(j\ge3).
\end{equation}
Set $q_d(m)=0$ when $(m,6d)>1$. Thus all these coefficients can be
computed using ordinary prime factorization and finite sums of Legendre
symbols. In particular, $q_d$ is not the ordinary M\"obius function.

Here is an explicit formula involving that function. Define
$\Lambda_d(r)=\prod_{p\mid r}\lambda_d(p)$ for squarefree $r$ coprime
to $6d$, with $\Lambda_d(1)=1$. For $(m,6d)=1$,
\begin{equation}\label{eq:integer-convolution}
 \boxed{\quad
 q_d(m)=\sum_{\substack{rt^2=m\\r\ \mathrm{squarefree}\\(r,t)=1}}
       \mu(r)\mu(t)^2\Lambda_d(r).
 \quad}
\end{equation}
At each prime, one chooses the constant, linear, or quadratic term of
the Euler polynomial. The restrictions in this formula make those choices
mutually exclusive.

Write $L(E_d,s)$ using the local factors of a minimal model, including at
primes dividing $6d$, and put
\[
 B_d(s)=\prod_{p\mid6d}L_p(E_d,s+\tfrac12).
\]
Then, for $\Re s>1$,
\begin{equation}\label{eq:reciprocal}
 \sum_{m\ge1}\frac{q_d(m)}{m^s}
 =\prod_{p\nmid6d}(1-\lambda_d(p)p^{-s}+p^{-2s})
 =\frac{B_d(s)}{L(E_d,s+\tfrac12)}.
\end{equation}
The finite factor $B_d(s)$ is holomorphic and nonzero for $\Re s>0$.
Deleting these primes therefore does not affect the zeros in question.
This remains true when the displayed equation for $E_d$ is not minimal.

Choose a complex-valued $W\in C_c^\infty((0,\infty))$, supported in a fixed
interval $[a,b]$ with $0<a<b$. The target sum is
\begin{equation}\label{eq:target}
 S_d(D)=\sum_{m\ge1}q_d(m)W(m/D).
\end{equation}
The weight restricts $m$ to size $D$. We want, for every $\varepsilon>0$,
$S_d(D)\ll_{d,W,\varepsilon}D^{11/12+\varepsilon}$.
Here $\ll$ allows a constant independent of $D$.

\subsection{The comparison with the zeta coefficients}
For zeta, the classical coefficients in the companion guide are
\[
 c(m)=\sum_{rt=m}\mu(r)\mu(t)\chi_{-3}(t),\qquad
 \chi_{-3}(t)=
 \begin{cases}0&3\mid t,\\1&t\equiv1\pmod3,\\-1&t\equiv2\pmod3.
 \end{cases}
\]
For $p\nmid6d$, the comparison is
\[
\begin{array}{c|cc|cc}
 &c(p)&c(p^2)&q_d(p)&q_d(p^2)\\\hline
 p\equiv1\pmod3&-2&1&-a_p(d)/\sqrt p&1\\
 p\equiv2\pmod3&0&-1&0&1
\end{array}
\]
The inert-prime sign really changes. For instance, for $d=1$,
$a_5(1)=0$, $a_7(1)=-4$, so
$q_1(5)=0$, $q_1(25)=1$, and $q_1(7)=4/\sqrt7$.
The character $\chi_{-3}$ records whether a prime splits; it does not
encode the point count of a particular elliptic curve.

\section{The triangular lattice and the curve's row}\label{sec:lattice}

Put $\omega=e^{2\pi i/3}$. An Eisenstein integer is an ordinary pair of
integers written as $n=A+B\omega$, with
\[
 \N(n)=|n|^2=A^2-AB+B^2,\qquad
 \alpha(n)=\frac{A+B\omega}{\sqrt{A^2-AB+B^2}}.
\]
The ring $\OO$ has unique factorization and six units
$\pm1,\pm\omega,\pm\omega^2$. Choose the unique \emph{primary} associate
of each element coprime to $3$:
\[
 n\equiv1\pmod3,\qquad\text{that is, }A\equiv1\pmod3,
 \quad B\equiv0\pmod3.
\]
Unless otherwise stated, column variables and their divisors are primary
and coprime to $6$; row variables run over all of $\OO$. Coprimality to
$2$ means that $A,B$ are not both even. A star on a column sum means
squarefree in $\OO$.

Define $\muk(n)$ by the usual factorization rule: it is zero if a prime
square divides $n$, and otherwise is $(-1)^r$ for a product of $r$ distinct
primary primes. For example,
\[
 \muk(-5)=-1\quad(\N(-5)=25),\qquad
 \muk(7)=1\quad(7=(1+3\omega)(-2-3\omega)).
\]
Thus $\muk(n)$ need not equal $\mu(\N(n))$.

For a prime $\pi\nmid6$, define $\chi_\pi(u)=(u/\pi)_6$ to be the sixth
root of unity congruent to $u^{(\N(\pi)-1)/6}$ modulo $\pi$, or zero if
$\pi\mid u$. Extend multiplicatively in the denominator, and write
$\chi_n(u)=(u/n)_6$. It is also multiplicative in $u$.
Powers of these characters retain their zero values, even when the
exponent is reduced to zero modulo $6$; $\chi_1(u)=1$ for all $u$.

Define the family
\begin{equation}\label{eq:Au}
 A_u(D)=\sum_n\muk(n)\alpha(n)\chi_n(u)W(\N(n)/D).
\end{equation}
It is the zeta family with one extra, explicitly given factor $n/|n|$.
For the desired curve, choose
\begin{equation}\label{eq:ud}
 u_d=(4d)^5,\qquad
 \psi_d(n)=\alpha(n)\chi_n(u_d)
          =\alpha(n)\overline{\chi_n(4d)}.
\end{equation}
We claim that $A_{u_d}(D)=S_d(D)$.

\paragraph{Where $u_d$ comes from.} At a split prime $p=\pi\bar\pi$, the point count
\eqref{eq:point-identity} says that $E_d$ contributes $\pi$ times the sixth root of unity
$\overline{\chi_\pi(4d)}$. In its Jacobi-sum proof the factor $d$ enters when the summation variable is
rescaled by $d$, and the factor $4$ enters through the discriminant $1-4a$ of $t^2-t+a$, which gives
$J(\kappa,\rho)=\kappa(4)J(\kappa,\kappa)$ for the cubic and quadratic characters $\kappa,\rho$ modulo $\pi$. The
family, however, twists by $\chi_n(u)$, with the row $u$ in the \emph{top} of the residue symbol. So we
need an Eisenstein integer $u$ with $\chi_n(u)=\overline{\chi_n(4d)}$ for every column $n$. Since
$\chi_n$ takes sixth roots of unity,
\[
 \overline{\chi_n(4d)}=\chi_n(4d)^5=\chi_n\bigl((4d)^5\bigr),
\]
and this is the choice $u_d=(4d)^5$. Multiplying a row by a sixth power changes it only at the
columns that share a prime with that sixth power: $\chi_n(uv^6)=\chi_n(u)\,\ind_{(n,v)=1}$. Such a change
alters finitely many Euler factors, so it does not move any zeros in $\Re s>0$. When $v$ is supported on
primes above $2$ and $3$, nothing changes at all, because the columns are coprime to $6$. For instance
$u=16d^5$ gives exactly the same row, since $(4d)^5=16d^5\cdot2^6$.
For $d=16$ this is $2^{24}=(2^4)^6$, a sixth power. So the curve $y^2=x^3+16$ (conductor $27$) is the row
$u=1$ itself, namely the plain sum $\sum_n\muk(n)\frac n{|n|}W(\N(n)/D)$, and $y^2=x^3+1$ is the row $u=16$.
More generally, every nonzero rational row $u$ is the row of a $j=0$ curve over $\mathbb Q$: since
$(4d)^5\equiv(4d)^{-1}$ modulo sixth powers, $u$ belongs to $E_d$ with $d=4^5u^5$. For rows that are not a
rational number times a sixth power, the coefficients $q_u(m)$ of~\eqref{eq:allrows} are in general not real.
Such a row is the $L$-function of a Hecke character of $\mathbb Q(\sqrt{-3})$, not of an elliptic curve over
$\mathbb Q$.

\subsection{Checking the claim prime by prime}
If $p\equiv1\pmod3$, then $p=\pi\bar\pi$, with both generators primary
and $\N(\pi)=p$. The classical CM point-count identity is
\begin{equation}\label{eq:point-identity}
 a_p(d)=\pi\overline{\chi_\pi(4d)}
       +\bar\pi\chi_\pi(4d)\qquad(p\nmid6d).
\end{equation}
This classical finite-field identity is proved with Jacobi sums; see
\cite[Chapter~18]{IR}, where primary means $\equiv2\pmod3$, which changes the sign of $\pi$.
We also checked it numerically for $480$ pairs $(d,p)$ (Section~\ref{sec:78}). Since
$\chi_{\bar\pi}(4d)=\overline{\chi_\pi(4d)}$, the two reciprocal
factors multiply to
\[
 (1-\psi_d(\pi)p^{-s})(1-\psi_d(\bar\pi)p^{-s})
 =1-\lambda_d(p)p^{-s}+p^{-2s}.
\]
If $p\equiv2\pmod3$, the primary prime is $-p$, of norm $p^2$.
For rational $x$ prime to $p$, $\chi_{-p}(x)=1$, because
$(p^2-1)/6=(p-1)(p+1)/6$. Since $\alpha(-p)=-1$, its factor is
$1+p^{-2s}$, as required by $a_p(d)=0$. If $p\mid d$, the symbols
vanish at primes above $p$, so their factors are deleted.

We have therefore checked the exact identities
\begin{equation}\label{eq:normgroup}
 q_d(m)=\sum_{\N(n)=m}\muk(n)\alpha(n)\chi_n(u_d),
 \qquad S_d(D)=A_{u_d}(D).
\end{equation}
The columns still run over primary $n$ coprime to $6$; the sum is empty
when $(m,6)>1$. The angular factor and sextic symbol are what distinguish
this identity from the zeta identity involving $\mu$ and $\chi_{-3}$.

\paragraph{The curve's coefficients are sextic residue symbols.} The identity~\eqref{eq:normgroup} for
the reciprocal has a companion for $L$ itself. Up to the finitely many Euler factors at $p\mid6d$,
$L(E_d,s+\frac12)=\prod_{\mathfrak p\nmid6}(1-\psi_d(\mathfrak p)\N(\mathfrak p)^{-s})^{-1}=\sum_n\psi_d(n)\N(n)^{-s}$,
so for $m$ coprime to $6d$,
\[
 \frac{a_m(E_d)}{\sqrt m}=\sum_{\N(n)=m}\alpha(n)\Bigl(\frac{u_d}n\Bigr)_6,\qquad
 q_d(m)=\sum_{\N(n)=m}\muk(n)\,\alpha(n)\Bigl(\frac{u_d}n\Bigr)_6 ,
\]
with $n$ primary and coprime to $6$. Here $a_m(E_d)$ are the Dirichlet coefficients of $L(E_d,s)$. The only
difference between the two is the M\"obius factor. At a prime the first identity is
$a_p(d)=\Tr\bigl(\pi\,(u_d/\pi)_6\bigr)$: the Frobenius eigenvalues are $\pi(u_d/\pi)_6$ and its conjugate, so the
sextic symbol selects which of the six rotations $\zeta\pi$ ($\zeta^6=1$) occurs. We checked the first identity
against the coefficients computed by PARI/GP for $d\in\{1,2,-3,7\}$ and all $m\le400$ coprime to $6d$.

Thus the coefficients of $L(E_d,s+\frac12)$ are exactly the sextic residue symbols of the zeta argument,
evaluated at the fixed row $u_d$ and rotated by the angle $n/|n|$. In the zeta argument the row $u$ is the
M\"obius sum $\sum_n\muk(n)(u/n)_6W(\N(n)/D)$ of the finite-order Hecke $L$-function attached to
$n\mapsto(u/n)_6$. This is the Artin $L$-function of a character of the Kummer extension
$\mathbb Q(\sqrt{-3},u^{1/6})/\mathbb Q(\sqrt{-3})$, one factor of the Dedekind zeta function of that
extension. The CM argument averages the same rows with every
coefficient multiplied by $n/|n|$, and its row $u_d$ is $L(s,\alpha\cdot(u_d/\cdot)_6)=L(E_d,s+\frac12)$, up to
finitely many Euler factors.


\section{The average, in coordinates and as a double sum}\label{sec:average}

The central average is
\begin{equation}\label{eq:M}
 M(D,H)=\sum_{0<\N(u)\le H}|A_u(D)|^2.
\end{equation}
There are about $H$ rows and about $D$ columns in each row.
Writing $Q(A,B)=A^2-AB+B^2$, this means exactly
\begin{equation}\label{eq:coordinates}
\begin{split}
 M(D,H)=\sum_{\substack{x,y\in\mathbb Z\\0<Q(x,y)\le H}}
 \Bigg|\sum_{\substack{A\equiv1\ (3),\ B\equiv0\ (3)\\
                          A,B\ \text{not both even}}}
 &\muk(A+B\omega)\frac{A+B\omega}{\sqrt{Q(A,B)}}\\[-1ex]
 &\times\left(\frac{x+y\omega}{A+B\omega}\right)_6
 W\left(\frac{Q(A,B)}D\right)\Bigg|^2.
\end{split}
\end{equation}
Thus it is a sum over four ordinary integer coordinates, using
factorization and exponentiation in finite fields.

Expanding the square, with the first column conjugated, gives
\begin{equation}\label{eq:double}
\begin{split}
 M(D,H)=\sum_{n_1,n_2}&\muk(n_1)\muk(n_2)
 \overline{\alpha(n_1)}\alpha(n_2)
 \overline{W(\N(n_1)/D)}W(\N(n_2)/D)\,S_H(n_1,n_2),\\
 S_H(n_1,n_2)&=\sum_{0<\N(u)\le H}
                 \overline{\chi_{n_1}(u)}\chi_{n_2}(u).
\end{split}
\end{equation}
The elliptic-curve factor in this double sum is simply
$\overline{n_1}n_2/(|n_1||n_2|)$. No abstract character notation is
needed to calculate it.

\subsection{An ordinary-integer description of every row}
If desired, put
\[
 q_u(m)=\sum_{\N(n)=m}\muk(n)\alpha(n)\chi_n(u),\qquad
 A_u(D)=\sum_{m\ge1}q_u(m)W(m/D).
\]
For fixed $u\ne0$, $q_u$ is multiplicative, vanishes on integers divisible
by $2$ or $3$, and for $p\ge5$ has the following local values:
\begin{equation}\label{eq:allrows}
\begin{array}{c|cc}
 &q_u(p)&q_u(p^2)\\\hline
 p=\pi\bar\pi\equiv1\pmod3
 &-\dfrac{\pi\chi_\pi(u)+\bar\pi\chi_{\bar\pi}(u)}{\sqrt p}
 &\chi_\pi(u)\chi_{\bar\pi}(u)\\[2ex]
 p\equiv2\pmod3&0&\chi_{-p}(u)
\end{array}
\end{equation}
All $q_u(p^j)$ with $j\ge3$ vanish. These formulas include the case
when a symbol is zero. At $u=u_d$ they reduce to~\eqref{eq:qd}.
For a general row, the symbols distinguish $\pi$ from $\bar\pi$;
ordinary $\mu$ and $\chi_{-3}$ alone do not determine them.

\subsection{The estimate to be supplied by analysis}
The analytic input is the following mean-square estimate: for every fixed
$0<\vartheta\le1/10$ and every $\varepsilon>0$,
\begin{equation}\label{eq:MS}
 \boxed{\quad M(D,D^{1+\vartheta})
       \ll_{W,\vartheta,\varepsilon}D^{2+\vartheta+\varepsilon}.
       \quad}
\end{equation}
The constant is controlled by finitely many derivatives of $W$.
The curve parameter does not occur in this estimate: all the curves are
selected by different rows of the same family.

The diagonal $n_1=n_2=n$ in~\eqref{eq:double} has
\[
 S_H(n,n)=\#\{u:0<\N(u)\le H,\ (u,n)=1\}\ll H.
\]
More precisely, elementary lattice counting gives
\[
 S_H(n,n)=\frac{2\pi}{\sqrt3}H
       \prod_{\pi\mid n}(1-\N(\pi)^{-1})
       +O\bigl(\tau(n)(\sqrt H+1)\bigr).
\]
Here $\tau(n)$ counts primary divisors. The factor $2\pi/\sqrt3$ is the
area of the disc divided by the lattice covolume. The diagonal is
$O_W(HD)$. Thus~\eqref{eq:MS} bounds the whole average on the diagonal
scale, up to a small power loss. It bounds the total off-diagonal
contribution, not the sum of its termwise absolute values.

\section{Sixth powers isolate the chosen elliptic curve}\label{sec:sixth}

Fix $u_0=u_d$ and $N_0=\N(u_0)=|4d|^{10}$. For any $v\ne0$,
\[
 \chi_n(u_0v^6)=\chi_n(u_0)\ind_{(n,v)=1}.
\]
Consequently the rows we need are \emph{the fixed row $u_0$ times sixth
powers}; they need not themselves be sixth powers.

Take $H=D^{1+\vartheta}$, put $Y=(H/N_0)^{1/6}$, and let
\[
 \mathcal P_Y=\{\pi\text{ primary prime}:\ \pi\nmid6d,
                                      \ Y/2<\N(\pi)\le Y\},
 \qquad J=|\mathcal P_Y|.
\]
For large $D$, the prime number theorem for $p\equiv1\pmod3$ gives
$J\asymp_d Y/\log Y$. The rows $u_0\pi^6$ are distinct and have norm
at most $H$. Their difference from the target is exactly
\[
 A_{u_0}(D)-A_{u_0\pi^6}(D)
 =\sum_{\pi\mid n}\muk(n)\alpha(n)\chi_n(u_0)W(\N(n)/D).
\]
There are $O_W(D/\N(\pi))$ possible multiples of $\pi$ in this sum,
because dividing by $\pi$ reduces the norm bound by $\N(\pi)$.
Every coefficient has absolute value at most one. Therefore
\begin{equation}\label{eq:approx}
 |A_{u_0}(D)-A_{u_0\pi^6}(D)|\ll_W D/Y.
\end{equation}
Using $|z|^2\le2|w|^2+2|z-w|^2$ for each selected row and the
nonnegativity of the mean square gives
\begin{equation}\label{eq:extraction}
 \boxed{\quad J|S_d(D)|^2
 \le2M(D,H)+O_W(JD^2/Y^2).\quad}
\end{equation}
There is no assumption that different rows cancel against each other.
Positivity is precisely what lets us extract the special rows.

Inserting~\eqref{eq:MS}, absorbing logarithms into $D^\varepsilon$, and
allowing constants depending on $d$, we get
\begin{align}
 |S_d(D)|^2
 &\ll_{d,W,\vartheta,\varepsilon}
          DH^{5/6}D^\varepsilon+D^2H^{-1/3}\notag\\
 &=D^{11/6+5\vartheta/6+\varepsilon}
                  +D^{5/3-\vartheta/3}.\label{eq:exponents}
\end{align}
Take square roots, then choose $\vartheta>0$ and the power loss small
enough. For every $\eta>0$ this proves
\begin{equation}\label{eq:Sbound}
 S_d(D)\ll_{d,W,\eta}D^{11/12+\eta}.
\end{equation}
\paragraph{Dependence on $d$.} The mean-square estimate~\eqref{eq:MS} does not involve $d$; the curve
enters only through the position of its row. Keeping track of $N_0=\N(u_0)$ in the steps above gives
\[
 S_d(D)\ll_{W,\vartheta,\varepsilon}N_0^{1/12}D^{11/12+5\vartheta/12+\varepsilon}+N_0^{1/6}D^{5/6-\vartheta/6},
 \qquad N_0^{1/12}=|4d|^{5/6},
\]
valid once $D^{1+\vartheta}\ge2^6N_0$, so that $Y\ge2$. This is nontrivial only for $D$ larger than about
$|d|^{10}$. That is a dependence on the chosen parameter $d$, not on the curve: $E_{dt^6}$ is isomorphic
to $E_d$ over $\mathbb Q$, so $|d|$ can be made arbitrarily large while the conductor stays fixed. If $d$ is
first replaced by its sixth-power-free part, then every prime $p\ge5$ dividing $d$ divides the
conductor $N(E_d)$ exactly twice, so $|d|\ll N(E_d)^{5/2}$, and $|d|^{10}$ is at most a fixed power of the
conductor. For the zero-free half-plane this dependence is harmless, because each $d$ is treated
separately.

The exponent is easy to remember: the average has size about $D^2$,
and about $D^{1/6}$ useful copies of the target occur in it, so the
target squared has size about $D^{2-1/6}=D^{11/6}$.

\subsection{Other perfect powers and the rational slice}
The exact identities
\[
 \chi_n(v^3)=(v/n)_2,\qquad \chi_n(v^2)=(v/n)_3
\]
identify cube and square rows as quadratic and cubic twists.
If $v$ is a nonzero rational integer, then
$(v/n)_2=(v/\N(n))$, an ordinary Jacobi symbol. Hence
\begin{equation}\label{eq:quadratic-twists}
 A_{u_dv^3}(D)=\sum_{m\ge1}q_d(m)
                       \left(\frac vm\right)W(m/D).
\end{equation}
The finite-field reason at an inert prime is that every nonzero element
of $\mathbb F_p$ is a square in $\mathbb F_{p^2}$; at a split prime it
is ordinary Euler's criterion.

The entire rational slice $u\in\mathbb Z$ has only $\asymp H^{1/2}$ rows,
since $\N(u)=u^2$, whereas the lattice family has $\asymp H$ rows.
Sixth powers of lattice elements number $\asymp H^{1/6}$; sixth powers
of rational integers number $\asymp H^{1/12}$. The primes used in
\eqref{eq:extraction} are lattice primes, not merely rational prime bases.
General rational numerators still require sextic residue symbols when
computing their coefficients at split primes.

\section{Why the bound excludes zeros}\label{sec:zeros}

\begin{proposition}
If~\eqref{eq:MS} holds for every smooth compactly supported complex
weight and every fixed $0<\vartheta\le1/10$, then
$L(E_d,s)\ne0$ for $\Re s>17/12$.
\end{proposition}
\begin{proof}
We have proved~\eqref{eq:Sbound}. Let
\[
 \widehat W(s)=\int_0^\infty W(y)y^s\,\frac{dy}{y},\qquad
 F_W(s)=\int_0^\infty S_d(D)D^{-s}\,\frac{dD}{D}.
\]
The second integral is holomorphic for $\Re s>11/12$. At the lower
endpoint there is no convergence problem: $S_d(D)=0$ for $D<1/b$.
For $\Re s>1$, absolute convergence and~\eqref{eq:reciprocal} give
\begin{equation}\label{eq:mellin-zero}
 L(E_d,s+\tfrac12)F_W(s)=\widehat W(s)B_d(s).
\end{equation}
Since $L(E_d,s)$ is entire, the identity theorem extends this equality
to $\Re s>11/12$.

Suppose $L(E_d,\rho)=0$ with $\Re\rho>17/12$, and put
$s_0=\rho-1/2$. Choose a nonzero nonnegative smooth function $\phi$
supported in $(1,2)$, and use $W(y)=y^{-s_0}\phi(y)$. Then
\[
 \widehat W(s_0)=\int_1^2\phi(y)\,\frac{dy}{y}>0.
\]
At $s=s_0$ the left side of~\eqref{eq:mellin-zero} is zero, while the
right side is nonzero. This contradiction proves the assertion.
\end{proof}
This explains both the allowance of complex weights and the shift by
$1/2$. The zero-free boundary asserted is strict; no conclusion at
$\Re s=17/12$ is being drawn.

\section{Poisson summation reveals the cubic Gauss sums}\label{sec:poisson}

We now explain the analytic work behind~\eqref{eq:MS}.
Choose a fixed smooth radial majorant $\Phi(\N(u)/H)$, nonnegative and
at least one on $\N(u)\le H$, rapidly decreasing at infinity, with
compactly supported planar Fourier transform. Put
\[
 e(z)=\exp(4\pi i\operatorname{Im}z/\sqrt3),\qquad
 \widehat\Phi(\N(w))=\frac2{\sqrt3}\int_{\mathbb C}
                         \Phi(|z|^2)e(-zw)\,dx\,dy.
\]
The smooth mean square over all $u\in\OO$, including zero, is an upper
bound for~\eqref{eq:M}.

Only squarefree columns occur. Write $n_i=gz_i$, where $g$ is their
primary greatest common divisor. Then $(z_1,z_2)=(g,z_1z_2)=1$, and
\[
 \overline{\chi_{n_1}(u)}\chi_{n_2}(u)
 =\ind_{(u,g)=1}\psi(u),\qquad
 \psi=\overline{\chi_{z_1}}\chi_{z_2}.
\]
This is a primitive character modulo $z_1z_2$. For a primitive character
modulo $q$, define the normalized finite Gauss sum
\[
 \gamma(\psi)=\frac1{\sqrt{\N(q)}}\sum_{x\bmod q}\psi(x)e(x/q).
\]
Inclusion--exclusion for $(u,g)=1$, followed by planar Poisson summation,
gives the exact formula
\begin{equation}\label{eq:poisson-exact}
\begin{split}
 &\sum_{u\in\OO}\Phi(\N(u)/H)\ind_{(u,g)=1}\psi(u)\\
 &\quad=\frac{H\gamma(\psi)}{\sqrt{\N(z_1z_2)}}
 \sum_{e_0\mid g}\frac{\muk(e_0)\psi(e_0)}{\N(e_0)}
 \sum_{h\in\OO}\overline{\psi(h)}
 \widehat\Phi\left(\frac{H\N(h)}{\N(e_0)\N(z_1z_2)}\right).
\end{split}
\end{equation}
The divisor $e_0$ must be kept. The zero frequency contributes only when
$z_1=z_2=1$; it gives the diagonal main term. There can also be nonzero
frequencies from these diagonal pairs, and they remain in the estimate.
The other frequencies satisfy
\[
 \N(h)\ll\frac{\N(e_0)\N(z_1z_2)}H
          \ll\frac{D^2}{H\N(g)}.
\]
This is the first shortening: $H>D$ produces frequencies of norm less
than about $D$.

\subsection{The arithmetic identity and the extra angular factor}
For squarefree $n$ define
\[
 \gamma_j(n)=\frac1{\sqrt{\N(n)}}
                 \sum_{x\bmod n}\chi_n(x)^j e(x/n),\qquad
 G(n)=\overline{\chi_n(4)}\gamma_3(n).
\]
In particular, $\gamma_2$ is a cubic Gauss sum. The finite-sum identity
used in both arguments is
\begin{equation}\label{eq:gauss}
 \muk(n)\gamma_{-1}(n)
 =\chi_n(-1)G(n)^{-1}\overline{\alpha(n)}\gamma_2(n).
\end{equation}
It combines the Jacobi-sum identity for a cubic Gauss sum with sextic
reciprocity. The M\"obius sign is essential here. The function $G$ has
absolute value one and depends only on the primary invertible residue
class modulo $72$.

Let $\mathfrak G$ be that finite group of classes. Define, quite
explicitly,
\[
 c_\xi=\frac1{|\mathfrak G|}\sum_{t\in\mathfrak G}
                         G(t^{-1})\overline{\xi(t)},\qquad
 a_\xi(n)=\overline{\alpha(n)}^{\,2}\gamma_2(n)\xi(n),
\]
where $\xi$ ranges over the characters of this fixed finite group.
For coprime squarefree $z_1,z_2$, the useful paired identity is
\begin{equation}\label{eq:paired}
\begin{split}
 &\muk(z_1)\muk(z_2)\overline{\alpha(z_1)}\alpha(z_2)
       \gamma(\overline{\chi_{z_1}}\chi_{z_2})\\
 &\hspace{3em}=\sum_\xi c_\xi a_\xi(z_1)\overline{a_\xi(z_2)}.
\end{split}
\end{equation}
This preserves the structure of a squared sum. The zeta version produces
$\overline\alpha\gamma_2$; the present version produces
$\overline\alpha^{\,2}\gamma_2$. The second factor of
$\overline\alpha$ comes from conjugating the first column of~\eqref{eq:Au}.

\subsection{The next average, with every variable displayed}
After resolving greatest common divisors, the averages have the form
\begin{equation}\label{eq:energy}
 \E(\mathcal H,X,F)=\frac1{XF}
 \ssum_{F\le\N(f)<2F}\ \sum_{0<\N(k)\le\mathcal H}
 \left|\ssum_n a_\xi(n)\chi_n(k)\chi_n(f)^4
                       W(\N(n)/X)\right|^2.
\end{equation}
Both $f,n$ are primary, squarefree, and coprime to $6$; $k$ runs over
all lattice points. The weight and the finite character $\xi$ have been
suppressed on the left. Put $\Sigma=XF$.

There are two sources of the fourth power. The divisor $e_0$ in
Poisson contributes $\chi_n(he_0^5)$. A second exclusion divisor $v_0$
uses the twisted multiplicativity
\[
 a_\xi(v_0n)=a_\xi(v_0)a_\xi(n)\chi_n(v_0)^4
 \qquad((v_0,n)=1).
\]
With $f=e_0v_0$ and $k=e_0h$, these combine as
$he_0^5v_0^4=kf^4$. Neither source can be omitted.

The initial scales are, for a further divisor-size parameter $B\ge1$,
\begin{equation}\label{eq:initial}
 X=\frac D{BF},\qquad \Sigma=\frac DB,\qquad
 \mathcal H\ll\frac{D^2}{HB^2}.
\end{equation}
The first Poisson reduction shows that the
estimate $\E\ll D^\varepsilon\Sigma$ in these ranges implies
\eqref{eq:MS}. Smooth kernels in two column variables are separated
using Fourier or Mellin integrals, with a controlled finite number of
derivatives; they are not simply discarded.

\section{The completed sums and the second theta derivative}\label{sec:theta}

For fixed $k,f$, put $\Psi(n)=\xi(n)\chi_n(k)\chi_n(f)^4$.
The completed sum is the following explicit double sum:
\begin{equation}\label{eq:T}
\begin{split}
 T(X;k,f)=\frac1{\sqrt X}\ssum_n\sum_b
 &\overline{\alpha(n)}^{\,2}\gamma_2(n)\Psi(n)
 \overline{\alpha(b)}^{\,6}\Psi(b)^3\N(b)^{1/2}\\
 &\hspace{3em}\times W\left(\frac{\N(n)\N(b)^3}X\right).
\end{split}
\end{equation}
Here $b$ is primary and coprime to $6$ but need not be squarefree.
All sums are finite. The term $b=1$ is exactly $X^{-1/2}$ times the
column sum in~\eqref{eq:energy}. The other terms are added because they
occur naturally in the theta expansion.

\subsection{Where those coefficients come from}
Put $\lambda=1+2\omega=i\sqrt3$ and
$\ce(z)=\exp(2\pi i(z+\bar z))$. Kubota's cubic theta function is a
Fourier series in a planar variable $z$ and a positive height $v$,
with terms $d_0(\ell)vK_{1/3}(4\pi|\ell|v)\ce(\ell z)$ in its
conjugate. Here $K_{1/3}$ is the usual modified Bessel function.
Patterson's coefficient formula gives
\begin{equation}\label{eq:theta-coeff}
 d_0(\lambda^{-3}nb^3)
  =3^{5/2}|b|\,\overline{(\lambda/n)_3}\,\gamma_2(n).
\end{equation}
Restricting $\lambda^3\ell$ to be primary and coprime to $6$ selects
exactly these indices. Other units and powers of $\lambda$ are excluded
from the twisted series by this restriction. In the reflected series, at
the three cusps, the indices lie in the larger set of
$\ell=u\lambda^mnb^3$ with $u$ a unit and $m\ge-4$. The coefficient tables of
Appendix~\ref{app:reflection} restrict them further; for example
$d_0(\lambda^{-4})=0$. The factor $(\lambda/n)_3$ depends only
on $n\bmod9$ and is canceled by inserting $(\lambda/\lambda^3\ell)_3$
in the twist. It is finite residue-class data, not a missing Gauss factor.

Differentiating $\ce(\ell z)$ twice with respect to $\bar z$ multiplies
it by $(2\pi i\bar\ell)^2$. Removing the resulting radial power by
a Mellin integral leaves $\overline{\alpha(\ell)}^{\,2}$.
Since $\ell=\lambda^{-3}nb^3$, its factors in $n,b$ are exactly
$\overline{\alpha(n)}^{\,2}\overline{\alpha(b)}^{\,6}$, up to the
fixed factor $\overline{\alpha(\lambda^{-3})}^{\,2}=-1$.
This explains every angular power in~\eqref{eq:T}.

\subsection{Why differentiating twice is legitimate}
The theta symmetry relates a rational translate $a/c+z$ to coordinates
\[
 z'=-\frac\delta c-\frac{\bar z}{c^2(v^2+z\bar z)},\qquad
 v'=\frac v{\N(c)(v^2+z\bar z)}.
\]
At $z=\bar z=0$, pure $\bar z$ derivatives of $\bar z'$ and $v'$
vanish, as do derivatives of $z'$ of order at least two, while
$\partial_{\bar z}z'=-(cv)^{-2}$. Therefore the ordinary chain rule gives
\begin{equation}\label{eq:chain}
 \left.\partial_{\bar z}^{\,2}F(z',\bar z',v')\right|_{z=\bar z=0}
   =(cv)^{-4}(\partial_{z'}^{\,2}F)
                  \left(-\delta/c,1/(\N(c)v)\right).
\end{equation}
No unwanted lower derivative survives. The derivative also kills any
constant Fourier term. Of the three expansions used, only the one at the
identity cusp has a nonzero constant term, and the derivative removes it.
This permits Mellin reflection without constant-term contributions.

For $V(y)=\sqrt y\,W(y)$ the reflected weight is
\begin{equation}\label{eq:weight}
 V^\sharp(x)=\frac1{2\pi i}\int_{\Re t=0}\widehat V(-t)
 \frac{\Gamma(4/3+t)\Gamma(5/3+t)}{\Gamma(4/3-t)\Gamma(5/3-t)}
 \left(\frac{(2\pi)^4x}{27}\right)^{-t}dt.
\end{equation}
This is the CM version of the gamma quotient in the zeta argument.
The reflected coefficients have angular factor $\alpha(\ell)^2$;
the fixed cusp prefactor is $-\overline{\alpha(c)}^{\,4}/81$, before
the finite Fourier coefficients and local phases are included.
The full reflected series and its coefficient tables are recorded in
Appendix~\ref{app:reflection}, so the main discussion can focus on the
lengths and the averaging.

\subsection{The local step from sextic to quadratic}
By reciprocity and finite congruence-class splitting, $\Psi$ is of the
form $\xi_0(n)\prod_{\pi\in\mathcal P}\chi_\pi(n)^{j_\pi}$.
The local reflected factors, apart from constant phases, are
\begin{equation}\label{eq:local}
 B_{\pi,j}(x)=
 \begin{cases}
  \chi_\pi(x)^{-j-2},&j\notin\{0,4\},\\
  \N(\pi)^{-1/2}\chi_\pi(x)^{-2},&j=0,\\
  \N(\pi)^{-1/2}(-1+\N(\pi)\ind_{\pi\mid x}),&j=4.
 \end{cases}
\end{equation}
In particular,
\[
 B_{\pi,1}(x)=\chi_\pi(x)^3=(x/\pi)_2.
\]
Away from the fixed factors of $6f$, squarefree row primes have exponent
$j=1$. These are the varying primes that become quadratic. The other
local factors are still present and have to be estimated.

If $Q=\N(c)$ is the norm of the reduced cusp denominator, the weight
in the reflected sum is $V^\sharp(\N(\ell)X/Q^2)$. Its effective norm
range is about $Q^2/X$, shorter than $X$ when $Q<X$.
Rapid decay gives a negligible tail when $X/Q^2$ is sufficiently large;
the inequality $Q^2<X$ alone does not specify a quantitative saving.

\section{The quadratic average and the completed bound}\label{sec:sieve}

The quadratic large sieve used here is
\begin{equation}\label{eq:qls}
 \ssum_{\substack{\N(k)\le K\\(k,6)=1}}
 \left|\ssum_{\N(n)\le U}\beta(n)\left(\frac nk\right)_2\right|^2
 \ll_\varepsilon(KU)^\varepsilon(K+U)\ssum_n|\beta(n)|^2.
\end{equation}
Both $k,n$ are primary and squarefree, the coefficients $\beta(n)$ are
arbitrary, and the column $n$ may be even. The squarefree restrictions
matter: square columns would give repeated values $1$ rather than
cancellation. This is the number-field quadratic large sieve of
\cite{GL}, in the form needed here.

Applying~\eqref{eq:qls} to the reflected sums, and retaining the factors
in~\eqref{eq:local}, gives the completed mean-square estimate
\begin{equation}\label{eq:completed}
 \boxed{\quad
 \sum_{0<\N(k)\le\mathcal H}|T(X;k,f)|^2
 \ll D^\varepsilon\left(\mathcal H+
                   \frac{\mathcal H^2\N(f)}X\right).
 \quad}
\end{equation}
All parameters are at least one and bounded by a fixed power of $D$;
a fixed finite smoothness norm of $W$ is understood. This is the CM form of the completed
mean-square estimate of~\cite{OAI1112}.

For clarity, passing from squarefree rows to all rows uses
$k=u s v_0^2$, with $u$ a unit and $s$ squarefree. The identity
\[
 \chi_n(v_0^2)\chi_n(f)^4=\chi_n(fv_0^2)^4
\]
absorbs the square into the auxiliary twist. The squarefree-row bound
allows that auxiliary argument to be nonsquarefree. Its contributions
are summed using $\sum_{v_0}\N(v_0)^{-2}<\infty$.

This is the main benefit of the theta transformation: the hard sextic
average has been reduced to an average of quadratic symbols. However,
\eqref{eq:completed} bounds the completed sum, whose cube terms still
have to be removed.

\section{Remove the cubes, then reduce the number of rows}\label{sec:descent}

One cannot discard the $b\ne1$ terms inside~\eqref{eq:T}: they could
cancel the term we want. Instead use the exact inversion
\begin{equation}\label{eq:inversion}
 \frac1{\sqrt X}\ssum_n a_\xi(n)\chi_n(k)\chi_n(f)^4W(\N(n)/X)
 =\sum_h\frac{\muk(h)\overline{\alpha(h)}^{\,6}\Psi(h)^3}{\N(h)}
                  T(X/\N(h)^3;k,f).
\end{equation}
Here $h$ is primary and coprime to $6$. To check the formula, substitute
\eqref{eq:T}, collect the total cube index $hb$, and use
$\sum_{h\mid b'}\muk(h)=\ind_{b'=1}$. The angular and character
factors are completely multiplicative, so this cancellation is exact.

Put
\[
 H_c^3=\min\left(X,\frac{X^2}{\mathcal H^2}\right),\qquad
 L_b=\frac X{\N(b)^3},\qquad\Sigma=XF.
\]
Assume $\mathcal H\le\Sigma$. For $\N(h)\le H_c$, the completed
bound at scale $X/\N(h)^3$ is at most
\[
 D^\varepsilon\left(\mathcal H+
              \frac{\mathcal H^2F\N(h)^3}X\right)
 \ll D^\varepsilon\Sigma.
\]
Cauchy--Schwarz and divisor sums absorb the sum over these small $h$.
Expanding the remaining terms gives the cube-removal inequality
\begin{equation}\label{eq:remove}
 \E(\mathcal H,X,F)\ll D^\varepsilon
 \left(\Sigma+
  \sup_{\substack{\N(b)>H_c\\L_b>1}}\E(\mathcal H,L_b,F)\right).
\end{equation}
An empty supremum is zero. Scales $L_b\le1$ are handled by counting;
if $\mathcal H^2\le X$, the supremum is empty.

\subsection{What two more Poisson summations do}
For a remaining scale $L=L_b$, smooth the row sum in the explicit average
\[
 \frac1{LF}\ssum_{F\le\N(f)<2F}\sum_{k\in\OO}
 \Phi(\N(k)/\mathcal H)
 \left|\ssum_n a_\xi(n)\chi_n(k)\chi_n(f)^4W(\N(n)/L)\right|^2.
\]
The reverse paired Gauss identity converts its coefficients back to
M\"obius coefficients $\muk(n)\overline{\alpha(n)}\xi_1(n)$.
For precision, the intermediate column sums, with their exclusions, are
\begin{equation}\label{eq:P}
 P_{C,e_0,t}(y;U)=\ssum_n\muk(n)\overline{\alpha(n)}\xi_1(n)
 \ind_{(n,t)=1}\overline{\chi_n(y)}\chi_n(C)^4\chi_n(e_0)
 U(\N(n)/\ell),\quad \ell=\frac L{\N(C)\N(t)}.
\end{equation}
Here $C,t$ are coprime squarefree primary elements and $e_0\mid C$.
After separation of smooth weights, the nonnegative average used is
\begin{equation}\label{eq:Q}
 \sum_{\substack{C,t\ \mathrm{squarefree}\\(C,t)=1}}
 \sum_{e_0\mid C}\frac{\mathcal H\N(C)}{\N(e_0)L^2F}
 \sum_{y\in\OO}\Phi(\N(y)/Y_{C,e_0})
                       |P_{C,e_0,t}(y;U)|^2,
\end{equation}
where the outer sum is restricted to $\N(Ct)\ll_W L$ and
\[
 Y_{C,e_0}=c_*\frac{\Sigma LF\N(e_0)}{\mathcal H\N(C)^2}.
\]
The constant $c_*$ is fixed and chosen to cover the support of the weights.

Here is the crucial enlargement. The first Poisson frequency $h$ has
size about $L^2/\mathcal H$, and combining it with $f$ as $y=hf^2$
gives a natural range about $L^2F^2/\mathcal H$. Suppressing the common
divisor factor $\N(e_0)/\N(C)^2$ and fixed constants, positivity permits
\begin{equation}\label{eq:enlarge}
 Y_{\rm natural}=\frac{L^2F^2}{\mathcal H}
 \quad\longrightarrow\quad
 Y_{\rm enlarged}=\frac{\Sigma LF}{\mathcal H},
 \qquad \frac{Y_{\rm enlarged}}{Y_{\rm natural}}=\frac{\Sigma}{LF}\ge1.
\end{equation}
A larger row range at this stage means a smaller frequency range after
the second Poisson summation. Without this enlargement, the two
summations need not give a strict reduction.

The second Poisson summation uses~\eqref{eq:paired} and returns to cubic
coefficients. Regrouping signed divisor contributions before taking
absolute values forces $f'\mid k'$. Locally the cancellation is
\[
 (1+\ind_{\pi\mid k'})-\ind_{\pi\mid k'}-\ind_{\pi\nmid k'}
 =\ind_{\pi\mid k'}.
\]
Thus the next row range is at least the size of its auxiliary divisor,
$\mathcal H'\ge F'\ge1$, and the column scales arising from removal of
exclusions also remain at least one. This keeps the next average in the
range to which the argument applies.

Precisely, the transfer estimate bounds the
average at $(\mathcal H,L,F)$ by $D^\varepsilon\Sigma$ times one plus
the supremum of normalized averages $\E(\mathcal H',X',F')/\Sigma'$,
where $\Sigma'=X'F'$ and
\begin{equation}\label{eq:transfer}
 \mathcal H'\le\frac{\mathcal HL}{\Sigma F},\qquad
 \frac{\mathcal H'}{\Sigma'}\le\frac{\mathcal H}{\Sigma},\qquad
 \Sigma'\le L.
\end{equation}
The finite character and the smooth weight may change, with controlled
derivatives. All three new parameters are at least one.

\subsection{Why this process terminates}
Suppose $\mathcal H\le\Sigma D^{-\kappa}$ for a fixed $\kappa>0$.
In a nonempty supremum in~\eqref{eq:remove},
$\N(b)^3>X^2/\mathcal H^2$. With $L=X/\N(b)^3$,
\begin{equation}\label{eq:contract}
 \mathcal H'\le\frac{\mathcal H}{\N(b)^3F^2}
 <\mathcal H\left(\frac{\mathcal H}{XF}\right)^2
 \le\mathcal H D^{-2\kappa}.
\end{equation}
The number of rows falls by a fixed power of $D$, while the ratio
$\mathcal H'/\Sigma'\le D^{-\kappa}$ is preserved.
Since the parameters start below a fixed power of $D$, finitely many
steps reach $\mathcal H\le1$. There are then at most six rows, and
counting gives $\E\ll X\le\Sigma$.
Distributing the allowed power loss among this bounded number of steps
proves
\[
 \E(\mathcal H,X,F)\ll D^\varepsilon\Sigma
 \qquad(\mathcal H\le\Sigma D^{-\kappa}).
\]
Finally, the initial ranges~\eqref{eq:initial} have
$\mathcal H/\Sigma\ll D^{-\vartheta}/B$.
For sufficiently large $D$, take $\kappa=\vartheta/2$.
This proves the required auxiliary estimate, hence the original mean
square~\eqref{eq:MS}, and hence the deduction in Section~\ref{sec:zeros}.

\paragraph{What changes from the zeta guide.}
The architecture is the same, but its angular factors change coherently:
\[
\begin{array}{c|cc}
 &\text{zeta guide}&\text{present CM guide}\\\hline
 \text{original coefficients}&\muk(n)&\muk(n)\alpha(n)\\
 \text{cubic coefficients}&\overline\alpha\gamma_2&\overline\alpha^{\,2}\gamma_2\\
 \text{theta derivative}&\partial_{\bar z}&\partial_{\bar z}^{\,2}\\
 \text{angular factor on cubes}&\overline{\alpha(b)}^3&\overline{\alpha(b)}^6\\
 \text{dual angular factor}&\alpha(\ell)&\alpha(\ell)^2
\end{array}
\]
These are genuine identities to check, not consequences of the factors
having absolute value one. Once checked, the same large-sieve and descent
estimates apply.

%======================== PART II: THE 7/8 BOUNDARY ========================
\clearpage
\section{Part II: the \texorpdfstring{$7/8$}{7/8} boundary for the same curves}\label{sec:78}

The remaining sections adapt OpenAI's $7/8$ argument~\cite{OAI78}, explained
for $\zeta$ in~\cite{G78}, to the curves $E_d$. The result is a wider
zero-free strip. Equation and line references to~\cite{OAI78} are to its
\TeX\ source. Throughout Part~II, $\mu=\muk$ is the M\"obius function on primary
elements, and $\chi_\bullet(u)$ denotes the character $n\mapsto\chi_n(u)$.

\begin{theorem}[claimed; see Section~\ref{sec:78status}]\label{thm:78}
Let $\eta$ be a finite-order Hecke character of $K=\mathbb Q(\sqrt{-3})$, and put
$\psi=\overline\alpha\,\eta$, that is, $\psi(n)=\overline{\alpha(n)}\,\eta(n)$ on
primary $n$ prime to the conductor of $\eta$. Then
\[
 L(s,\psi)\ne0\qquad(\Re s>7/8).
\]
In particular, for every nonzero integer $d$,
\[
 L(E_d,s)\ne0\qquad(\Re s>11/8),
\]
a zero-free strip $11/8<\Re s<3/2$ along the right edge of the critical strip.
\end{theorem}

\paragraph{How the curves enter.} Put $\psi_d^*(n)=\overline{\alpha(n)}\chi_n(4d)$.
If $p\nmid6d$ splits as $p=\pi\bar\pi$ with $\pi$ primary, then
$\sqrt p\,\bigl(\psi_d^*(\pi)+\psi_d^*(\bar\pi)\bigr)=a_p(d)$; if $p\equiv2\pmod3$,
then $\psi_d^*(-p)=-1$, matching $a_p(d)=0$. So $L(E_d,s+\frac12)=L(s,\psi_d^*)$
up to finitely many Euler factors, which do not vanish in $\Re s>0$.
(We checked both statements for $480$ split pairs $(d,p)$, testing both primes above each $p$, and
$112$ inert pairs.)
The character $\psi_d^*$ is the complex conjugate of the character
$\psi_d=\alpha\,\overline{\chi_\bullet(4d)}=\alpha\,\chi_\bullet\bigl((4d)^5\bigr)$ used in
Sections~\ref{sec:lattice}--\ref{sec:descent}. The two have the same
$L$-function: since $4d$ is rational, $\psi_d^*(\bar\pi)=\overline{\psi_d^*(\pi)}$, so both
characters give the same Euler factor at each split prime, and $\psi_d^*(-p)=-1$ is real.
More generally, every unitary Hecke character of $K$ whose infinity type is
$(\bar z/|z|)^{\pm1}$ is of the form $\overline\alpha\eta$ or its conjugate
$\alpha\bar\eta$. So the theorem covers every weight-two newform with complex
multiplication by $K$. In particular it covers every elliptic curve over
$\mathbb Q$ with complex multiplication by an order of $K$ (the $j$-invariants
$0$, $54000$ and $-12288000$).

\paragraph{What ``proof'' means in Part II.} We do not rewrite the
$16{,}700$ lines of~\cite{OAI78}. Instead we do three things:
\begin{enumerate}
\item we prove the identities in which the angular factor changes a
 computation, and test them numerically;
\item we state each ingredient of~\cite{OAI78} in the form it takes for a CM
 target;
\item we checked, section by section, that each proof in~\cite{OAI78} uses only
 properties that CM characters share with finite-order characters, or we record
 the change.
\end{enumerate}
No step was found that needs a new idea. Section~\ref{sec:78status} lists the
statements that still have to be written out in full, at the generality
of~\cite{OAI78}.

\subsection{The family and the number \texorpdfstring{$\beta_*$}{beta*}}
Let $\mathcal F_0$ be the set of primitive finite-order Hecke characters of $K$, and
let $\mathcal F_{\rm CM}$ be the set of primitive Hecke characters of $K$ with infinity type
$\bar z/|z|$, that is, the primitive characters inducing $\overline\alpha\eta$ with $\eta$ of finite
order. Put
\[
 \beta_*=\sup\bigl(\{1/2\}\cup\{\Re\rho:\ L(\rho,\psi)=0,\ 1/2\le\Re\rho\le1,\
 \psi\in\mathcal F_{\rm CM}\}\bigr).
\]
Three properties make this family suitable.
\begin{itemize}
\item \emph{Closure under finite-order twists.} If $\psi\in\mathcal F_{\rm CM}$ and $\theta$
 has finite order, then the primitive character inducing $\psi\theta$ lies in
 $\mathcal F_{\rm CM}$. Every denominator row below is such a twist.
\item \emph{Conjugation.} $\overline\psi=\alpha\bar\eta$ and
 $L(s,\overline\psi)=\overline{L(\bar s,\psi)}$, so conjugates contribute no new real parts
 of zeros.
\item \emph{No poles.} Every $L(s,\psi)$ with $\psi\in\mathcal F_{\rm CM}$ is entire. No
 member is principal, so the principal-character clauses of~\cite{OAI78} disappear.
\end{itemize}
The rows of the argument also contain finite-order $L$-functions. For them we do
not use $\beta_*$. Instead we use OpenAI's theorem~\cite[Theorem~thm:main]{OAI78}:
they have no zeros in $\Re s>7/8$. So, under the contradiction hypothesis
$\beta_*>7/8$, every $L$-function that occurs has its zeros in $\Re s\le\beta_*$.
In~\cite{OAI78} this was automatic from the definition of $\beta_*$. Here it is one of two
places where the finite-order theorem is an input; the other is the zero-free hypothesis for the
slot characters in the plain moment (Section~\ref{sec:78moments}).

The continuation criterion~\cite[Prop.~lem:continuation-criterion]{OAI78}
(see~\cite[Proposition~1]{G78}) holds verbatim with $\eta$ replaced by $\psi$. Its proof uses
only three things:
\begin{itemize}
\item the definition of $\beta_*$;
\item the bound $|1/L^S(s,\psi)|\le\zeta_K(2)$ on $\Re s\ge2$;
\item the nonvanishing of finitely many Euler factors in $\Re s>0$.
\end{itemize}
It therefore suffices, for each $\psi\in\mathcal F_{\rm CM}$, to construct a function
$J_\psi(Z)$ with the low and high estimates of the criterion, for $\sigma_0=7/8$ and
$C(s)=s-11/16$, with margins independent of $\psi$.

The construction is that of~\cite{OAI78}, with three changes:
\begin{itemize}
\item the weight of the target (Section~\ref{sec:78probe});
\item the bookkeeping of the characters attached to rows
 (Section~\ref{sec:78rows});
\item an angular factor in the moment estimates
 (Section~\ref{sec:78moments}).
\end{itemize}
Section~\ref{sec:78final} checks that the final inequality is unchanged.

\section{The probe for a CM target}\label{sec:78probe}

\subsection{The weight}
In~\cite{OAI78} the target $\eta$ enters the probe~\cite[eq.~general-probe]{OAI78},
\cite[\S3.1]{G78}, only through the weight on the completed index $A=cn^3$:
$\gamma_2(c)\overline{\alpha(A)}\eta(A)$. For the target $\psi=\overline\alpha\eta$ the weight is
therefore
\begin{equation}\label{eq:cm-weight}
 \gamma_2(c)\overline{\alpha(cn^3)}\,\psi(cn^3)
 =\gamma_2(c)\overline{\alpha(c)}^{\,2}\,\overline{\alpha(n)}^{\,6}\,\eta(c)\eta(n)^3 .
\end{equation}
Everything else in the probe is unchanged:
\begin{itemize}
\item the additive polynomial with its sextic Gauss sums;
\item the averaging variable $m$;
\item the fixed phases;
\item the scales $X,Y,Z$.
\end{itemize}

\paragraph{Why $\overline\alpha\eta$ and not $\alpha\bar\eta$.} The two targets have conjugate
zeros, so either could be used. With $\alpha\bar\eta$ the weight would be
$\gamma_2(c)\bar\eta(cn^3)$, with no angular factor. The completed sums would then be
sums of coefficients of the cubic theta function itself. Its expansion at the identity cusp has a
nonzero constant term; the expansions at the other two cusps have none. In~\cite{OAI78} a
$\bar z$-derivative removes that term (L2331). Without a derivative it enters the reflection of every
translate attached to the identity cusp, and produces a main term that nothing in~\cite{OAI78} controls.
With the target $\overline\alpha\eta$ the weight carries $\overline\alpha^{\,2}$. The completed sums are then those of
$\partial_{\bar z}^{\,2}\theta$, which is the input of Section~\ref{sec:theta}, and the constant term is
removed.

\subsection{Second evaluation: the rows}
\begin{proposition}\label{prop:cm-rows}
With the weight~\eqref{eq:cm-weight}, Poisson summation in $m$ gives the
representation~\cite[eq.~probe-poisson-identity]{OAI78} of the probe, with row $u$
equal to
\[
 \frac{\zeta_K^S(6z)\,L^S\bigl(w,\chi_\bullet(u)\bigr)}
      {L^S\bigl(x,\psi\overline{\chi_\bullet(u)}\bigr)}\,\mathcal H_{\psi,u}(x,w,z),
\]
where $\mathcal H_{\psi,u}$ is the Euler product of~\cite[Lemma~lem:local-euler]{OAI78}
with $\eta(p)$ replaced by $\psi(p)$. It has the following properties:
\begin{itemize}
\item it converges in the same regions;
\item it satisfies $\mathcal H_{\psi,u}\ll\N(u)^\varepsilon$ there;
\item $\mathcal H_{\psi,1}=1+O(P_0^{-4/5})$ in the second region of~\cite{OAI78} (L4056).
\end{itemize}
\end{proposition}
\begin{proof}
\cite{OAI78} evaluates the coefficient prime by prime. The target enters only through
the unimodular numbers $\eta(p)$ and complete multiplicativity, and no ray-class
condition on $p$ is used (L3871).

The angular factor of the weight is absorbed by the Gauss-sum identity
\begin{equation}\label{eq:signal-cm}
 \overline{\alpha(c)}\,\gamma_2(c)\,\gamma_1(c)=\mu(c)\,G(c),
\end{equation}
a rearrangement of the second identity $\gamma_1\gamma_2=\mu\alpha G$ in~\cite[eq.~signal]{OAI78}.
Here $\gamma_1$ is the normalized sextic Gauss sum and $G$ a fixed finite phase. For the
weight~\eqref{eq:cm-weight}, $\overline{\alpha(c)}^{\,2}\gamma_2(c)\eta(c)\gamma_1(c)=\mu(c)G(c)\psi(c)$.
So every entry of the local table of~\cite{OAI78} holds with $\eta(p)$ replaced by
$\psi(p)$:
\begin{itemize}
\item The squarefree part gives $-\psi(p)\overline{\chi_p(u)}\N(p)^{-x}$, which is the
 Euler factor of $1/L(x,\psi\overline{\chi_\bullet(u)})$.
\item The cube part gives the phases $a_p=\overline{\alpha(p)}^{\,3}\psi(p)^3$ and
 $b_p=\overline{\alpha(p)}^{\,2}\psi(p)^2\overline{\chi_p(4)}$.
\item The relation $b_p^3=a_p^2$ used there still holds, since $\chi_p(4)^3=(4/p)_2=1$.
\end{itemize}
The convergence regions and the bounds for $\mathcal H$ use only $|\psi(p)|=1$.
\end{proof}
We checked~\eqref{eq:signal-cm} numerically at $60$ primes (error below
$4\cdot10^{-15}$). The proposition gives the following.
\begin{itemize}
\item The numerators $L(w,\chi_\bullet(u))$ have finite order. They are the same as for
 $\zeta$.
\item The denominators $L(x,\psi\overline{\chi_\bullet(u)})$ belong to $\mathcal F_{\rm CM}$ and
 are never principal.
\item The principal row $u=1$ is $\zeta_K^S(w)\zeta_K^S(6z)/L^S(x,\psi)$. Its residues at
 $w=1$ and $z=1/6$ give the signal with $H_\psi(x)/L^S(x,\psi)$, and with the same
 positive constant $c_S$. The positivity $M(1/6)>0$ is a property of the averaging
 weight alone.
\end{itemize}

\subsection{First evaluation: completed rows with two derivatives}
Cauchy--Schwarz splits the probe \cite[\S3.2]{G78} into two mean squares. The first is
the additive one, which does not involve the target. The second is the mean square of
the completed rows
\[
 B_m=\sum_{\substack{c\ {\rm sf}\\n}}
 \frac{\gamma_2(c)\,\overline{\alpha(cn^3)}^{\,2}\,\nu(cn^3)\,\chi_{cn^3}(m)}{\sqrt{\N(c)}\,\N(n)}
 \,V\Bigl(\frac{\N(c)\N(n)^3}Z\Bigr).
\]
Here $\nu$ is a finite-order character whose conductor is supported on the fixed set
$S$; it is the finite part $\eta$ times the fixed phases. The angular factor must stay in
the weight, not in $\nu$: the theta step of~\cite{OAI78} expands $\nu$ over residue
classes, which is impossible for $\overline\alpha$.

\begin{proposition}\label{prop:cm-completed}
$\sum_{\N(m)\le Q}|B_m|^2\ll Z^{\varepsilon}(Q+Q^2/Z)$, exactly as
in~\cite[Lemma~lem:unmarked-completed-moment]{OAI78}.
\end{proposition}
The compensated probe needs more than this. Its marked rows involve sums over tuples of slot primes,
divisibility conditions on the completed index, and a scale that depends on the tuple. In~\cite{OAI78}
their bound comes from separate lemmas: lem:hybrid-energy, lem:reflected-energy and
lem:probe-row-norm, through Cor.~cor:marked-reflection. Their CM versions, uniform in the tuples, the
masks and the smooth weights, are item~1 of Section~\ref{sec:78status}. They have not been written out
here.


\cite{OAI78} proves this through its completed reflection
\cite[Prop.~prop:completed-reflection]{OAI78}, which is the transformation formula for
$\partial_{\bar z}\theta$. For~\eqref{eq:cm-weight} one uses instead $\partial_{\bar z}^{\,2}\theta$,
as in Section~\ref{sec:theta} and Appendix~\ref{app:reflection}. The changes are:

\begin{center}
\begin{tabular}{lll}
\toprule
 & $\zeta$ and finite order~\cite{OAI78} & CM target\\\midrule
derivative & $\partial_{\bar z}\theta$ & $\partial_{\bar z}^{\,2}\theta$\\
direct Mellin factor & $\Gamma(s+\frac13)\Gamma(s+\frac23)$ & $\Gamma(s+\frac56)\Gamma(s+\frac76)$\\
reflected kernel & $\prod_\pm\frac{\Gamma(1+t\pm1/6)}{\Gamma(1-t\pm1/6)}$
 & $\frac{\Gamma(4/3+t)\Gamma(5/3+t)}{\Gamma(4/3-t)\Gamma(5/3-t)}$\\
cusp scalar & $-\frac{i}{81}\overline{\alpha(c)}^{\,2}$ & $-\frac1{81}\overline{\alpha(c)}^{\,4}$\\
angular factor, original side & $\overline{\alpha(n)}\,\overline{\alpha(b)}^{\,3}$ & $\overline{\alpha(n)}^{\,2}\overline{\alpha(b)}^{\,6}$\\
angular factor, dual side & $\alpha(\mu)$ & $\alpha(\mu)^2$\\\bottomrule
\end{tabular}
\end{center}

Both kernels satisfy $|R(\sigma+iu)|\ll_A(1+|u|)^{4\sigma}$ on $-1/4\le\sigma\le A$. The CM kernel is
holomorphic for $\Re t>-4/3$; OAI78's is holomorphic for $\Re t>-5/6$. So the reflected weight obeys
the same bound $\ll\min(x^{1/4},x^{-A})$, with the same seminorms. The other ingredients
carry over as follows.
\begin{itemize}
\item \emph{Dual lengths.} They are unchanged.
\item \emph{Quadratic large sieve.} The dual angular factor $\alpha(\mu)^2$ is a unimodular
 coefficient of the dual index, so it is absorbed into the arbitrary coefficients of the
 sieve.
\item \emph{Marks.} Slot primes in the completed index are handled in~\cite{OAI78} through
 $\ind_{q\mid A}=1-\chi_q(A)^0$, which never requires splitting $\overline\alpha$.
\item \emph{Phases.} The extra $\overline{\alpha(c)}^{\,4}$ is multiplicative over primes, so the
 row and tuple phases still separate.
\end{itemize}
The balanced low estimate $|I_\psi(Z^{1/2},Z^{1/2},Z)|\ll Z^{1/4+\varepsilon}$ follows from the
proposition as in~\cite{OAI78}. The compensated low estimate $|I_{\rm mod}(Z)|\ll Z^{3/16+\varepsilon}$
follows in the same way once the marked lemmas above are established in CM form.
We checked numerically:
\begin{itemize}
\item the Bessel--Mellin identities for one and two derivatives (relative error
 $3\cdot10^{-26}$);
\item the two reflection constants;
\item the Stirling exponents of the kernels.
\end{itemize}

\section{Rows, zeros and prime slots}\label{sec:78rows}

\subsection{Growth of the CM \texorpdfstring{$L$}{L}-functions}
\begin{lemma}\label{lem:cm-growth}
Let $\psi\in\mathcal F_{\rm CM}$ have conductor $\mathfrak f$, and put $Q=\N\mathfrak f$. Then
\[
 \Lambda(s,\psi)=(3Q)^{s/2}(2\pi)^{-s}\,\Gamma(s+\tfrac12)\,L(s,\psi)
\]
is entire and $\Lambda(s,\psi)=\epsilon(\psi)\Lambda(1-s,\overline\psi)$ with $|\epsilon(\psi)|=1$. Moreover
$|L(\sigma+it,\psi)|\ll Q^{3/5}(3+|t|)^2$ for $-1/10\le\sigma\le11/10$, which is the bound
of~\cite[Lemma~lem:hecke-strip-growth]{OAI78} for finite-order characters.
\end{lemma}
\begin{proof}[Sketch]
The functional equation is Hecke's, for characters of infinity type
$\bar z/|z|$~\cite[Ch.~VII, \S8]{Neukirch}.

\emph{Entire.} $L(s,\psi)$ is the Mellin transform of a theta series
$\sum_{a\in\OO}\phi(a)\,\bar a\,e^{-cv\N(a)}$, where $c>0$ and $\phi$ is the finite part, extended
to elements so that $\phi(\zeta a)\bar\zeta=\phi(a)$ for units $\zeta$. Poisson summation takes
$\bar z\,e^{-cv|z|^2}$ to a constant multiple of $\bar\xi\,e^{-c'|\xi|^2/v}v^{-2}$, which vanishes
at $\xi=0$. So there is no constant term, and hence no pole.

\emph{Growth.} On $\Re s=-1/10$ the gamma quotient satisfies
$|\Gamma(3/2-s)/\Gamma(1/2+s)|\asymp(3+|t|)^{6/5}$. That is the exponent of $|\Gamma(1-s)/\Gamma(s)|$ for
finite-order characters, so the functional equation gives
$|L(-1/10+it,\psi)|\ll Q^{3/5}(3+|t|)^{6/5}$, as in~\cite{OAI78}. On $\Re s=11/10$ the Euler product is
bounded. The damped Phragm\'en--Lindel\"of argument of~\cite{OAI78}, applied to
$Q^{-3/5}L(s,\psi)/(s+2)^2$, then gives the strip bound.
\end{proof}
Consequently the logarithmic control of~\cite[Lemma~lem:logarithmic-control]{OAI78}
holds for $\mathcal F_{\rm CM}$, with the same conductor exponents and without its
principal clauses. Only fixed polynomial orders in the height change, and these may
depend on the target.

\subsection{The characters attached to a row}
In~\cite{OAI78} each row $u$ carries a finite set of characters: the twists
$\theta\chi_\bullet(u)^{\pm1}$ by members $\theta$ of the finite group
$\Theta_0=\langle\eta,\widehat T\rangle$, where $\widehat T$ consists of fixed characters
supported over $2$ and $3$. A row is \emph{bad} if one of these has a zero of bounded
height well to the right of $1/2$. For a CM target the group generated by $\psi$ is
infinite. We use instead the finite set
\[
 \mathcal X_u=\bigl\{\theta\chi_\bullet(u)^{\pm1}\bigr\}\cup
 \bigl\{\overline\alpha\,\theta\,\overline{\chi_\bullet(u)},\ \alpha\,\theta\,\chi_\bullet(u)\bigr\},
 \qquad\theta\in\Theta_0=\langle\eta,\widehat T\rangle .
\]
It has three properties.
\begin{itemize}
\item It contains the numerator, the denominator $\psi\overline{\chi_\bullet(u)}$, and their
 conjugates.
\item It is finite and closed under conjugation. Together with Lemma~\ref{lem:cm-growth}, the
 conductor bound $Q_\psi\ll\N(u)$, and the fact that no CM member is principal, these are all
 that the detector of~\cite[Section~sec:detector]{OAI78} uses.
\item Its finite-order members have no zeros in $\Re s>7/8$, and its CM members none in
 $\Re s>\beta_*$. This gives the bin ceiling $a\le\beta_*$ of \cite[\S3.3]{G78}.
\end{itemize}
A zero of a CM member produces the two witnesses of \cite[\S3.3]{G78}, now with the
angular factor:
\[
 M_r=U^{-r/2}\sum_n\mu(n)\overline{\alpha(n)}\theta(n)\overline{\chi_n(u)}\,W\Bigl(\frac{\N(n)}{U^r}\Bigr),
 \qquad
 S_m=U^{-m/2}\sum_l\overline{\alpha(l)}\theta(l)\overline{\chi_l(u)}\,W'\Bigl(\frac{\N(l)}{U^m}\Bigr),
\]
or their conjugates. In each, $\overline\alpha$ appears together with $\overline\chi$, or $\alpha$ together
with $\chi$. We call such sums \emph{co-oriented}; this matters in
Section~\ref{sec:78moments}.

\subsection{The compensated probe and the slots}
The compensated probe \cite[\S4.2]{G78} is formed with $\overline{\psi(p_{J^c})}$ in place of
$\overline{\eta(p_{J^c})}$. In row $u$ a slot prime $p\nmid u$ contributes the local factor
$\psi(p)\overline{\chi_p(u)}\N(p)^{-x}$ of the denominator, and $\overline{\psi(p)}\psi(p)=1$. So the
marked-minus-rescaled difference is $-\overline{\chi_p(u)}$ plus small terms, exactly as
for~$\zeta$. The slot amplitudes
\[
 Q_i(u)=P_i^{-1/2}\sum_{p\in\text{slot }i}\Omega_i\Bigl(\frac{\N(p)}{P_i}\Bigr)
 \Bigl(\frac{\N(p)}{P_i}\Bigr)^{z-1}\overline{\Bigl(\frac up\Bigr)_6}
\]
are the same finite-order prime polynomials as in \cite[\S4.4]{G78}. Two consequences
follow.
\begin{enumerate}
\item The principal normalizer $A_T(Z)$ is unchanged.
\item The pointwise bound $|Q_i(u)|\ll P_i^{\delta/2+o(1)}$ of~\cite[Lemma~lem:bin-prime-bound]{OAI78}
 is unchanged, since the slot characters $\theta\overline{\chi_\bullet(u)}$ lie in $\mathcal X_u$.
 Separately, the uniform slot bound inside the plain moment (\cite[eq.~old-eq:3.5]{OAI78}) needs a
 zero-free region only for finite-order characters. OpenAI's theorem supplies it at $7/8$, so
 $\kappa=3/4$ may be used there and in the plain capacity, in place of $\kappa=2\beta_*-1$.
\end{enumerate}

\section{Counting bad rows with an angular factor}\label{sec:78moments}
\cite{OAI78} counts bad rows with the sextic large sieve of \cite[\S3.3]{G78} and the three averages of \cite[\S5.1]{G78}. For rows whose zero comes from a
finite-order member they apply unchanged. For CM members they are needed with the angular
factor on the summation variable, but not on the slot primes.

\paragraph{1. The sextic large sieve}~\cite[Lemma~lem:sextic-large-sieve]{OAI78}.
It allows arbitrary coefficients, so it is unchanged. This is the only row count used in the $11/12$ stage of~\cite{OAI78} (its Part~I).

\paragraph{2. The inverse moment with prime factors.} For either sign $\epsilon=\pm1$, and under the
conditions $r+2z\le m-c_1$, $2r+8z\le3m-c_2$ of~\cite[Lemma~lem:marked]{OAI78}, put
\begin{align*}
 M_u&=Z^{-r/2}\sum_n\mu(n)\,\alpha(n)^{\epsilon}\nu(n)\chi_n(u)^{\epsilon}
 W\Bigl(\frac{\N(n)}{Z^r}\Bigr),\\
 Q_u&=\prod_{i}Z^{-z_i/2}\sum_{p}a_i(p)\nu(p)\chi_p(u)^{\epsilon}W_i\Bigl(\frac{\N(p)}{Z^{z_i}}\Bigr),
\end{align*}
where $z=\sum_iz_i$, the slot sets $\mathcal P_i$ are disjoint and avoid $S$, $\N(p)\asymp Z^{z_i}$ for
$p\in\mathcal P_i$, and the coefficients $a_i(p)$ are independent of $u$ with $|a_i(p)|\le1$. Then
\begin{equation}\label{eq:cm-inverse}
 \sum_{0<\N(u)\ll Z^m}|M_uQ_u|^2\ll Z^{m+\varepsilon}.
\end{equation}
\emph{Why.} The proof in~\cite{OAI78} alternates between two coefficient classes: Möbius coefficients
$\mu\nu$, and cubic coefficients $\overline\alpha\gamma_2\nu$ attached to the theta function. Each
switch uses~\eqref{eq:signal-cm} or its companion
\[
 \mu(n)\gamma_{-1}(n)=\chi_n(-1)G(n)^{-1}\,\overline{\alpha(n)}\,\gamma_2(n),
\]
which we checked at $60$ primes. \cite{OAI78} runs the orientation $\mu\nu$ with $\overline\chi$. With the
co-oriented factor $\overline\alpha$ on the Möbius variable ($\epsilon=-1$) the classes become
$\mu\overline\alpha\leftrightarrow\overline\alpha^{\,2}\gamma_2$; for $\epsilon=+1$ they are the complex conjugates. These are the dual coefficient and the cube
factor $\overline\alpha^{\,6}$ of Sections~\ref{sec:poisson}--\ref{sec:descent}, and the theta steps use
$\partial_{\bar z}^{\,2}\theta$ as in Proposition~\ref{prop:cm-completed}.

A slot prime that is merged with the Möbius variable multiplies its coefficient by a sign and the
unimodular factor $\overline{\alpha(p)^{\epsilon}}$. That coefficient is still bounded, independent
of $u$, and of product form, which is all~\cite{OAI78} uses. The descent, its masks and punctures,
and the propagation of seminorms never see $\alpha$.

\emph{Orientation.} The anti-oriented sum, with $\alpha$ paired with $\overline\chi$, would turn into
plain $\gamma_2$, since $\alpha\cdot\overline\alpha\gamma_2=\gamma_2$. That is the theta function without a
derivative, with its constant terms. So the clause ``for either sign'' of~\cite{OAI78} does not extend
to that case. It is never needed, because all CM witnesses are co-oriented.

\paragraph{3. Sixth-power amplification}~\cite[Lemma~lem:inverse-amplification]{OAI78}.
This is unchanged. The identity $\chi_n(ua^6)=\chi_n(u)$ for $(a,n)=1$ does not involve $\alpha$. Its
input is~\eqref{eq:cm-inverse} without slots; for $\epsilon=1$ and trivial $\nu$, the mean square~\eqref{eq:MS} of Part~I of these notes is the
special case of row range at most $D^{11/10}$.

\paragraph{4. The plain fourth moment with prime factors}~\cite[Lemma~lem:plain]{OAI78}.
The CM statement needed is the following. Fix $e=\pm1$ and $3/4\le\kappa\le1$. For rows $k$ with
$0<\N(k)\ll Z^m$, write $\psi_k=\tau\chi_\bullet(k)$ for the finite-order row character, and put
\[
 S_k(n_j)=Z^{-n_j/2}\sum_l\alpha(l)^e\psi_k(l)W_j\Bigl(\frac{\N(l)}{Z^{n_j}}\Bigr)\qquad(j=1,2).
\]
Let $Q$ be a product of slot polynomials in $\psi_k$ with no angular factor, of total length $z$, whose
coefficients are fixed combinations of characters in $\Theta_0$. Let $\mathcal R_z$ be the rows of
\cite{OAI78}, defined through the finite-order part $\psi_k$. Then
\[
 \sum_{k\in\mathcal R_z}|S_k(n_1)S_k(n_2)Q|^2\ll Z^{M+\varepsilon},\qquad M=m+q,
\]
with the same hypotheses and uniformity as in~\cite{OAI78}: no restriction when $z=0$;
$n_1+n_2+6\kappa z\le M$ and slots of length at most the mesh when $z>0$; uniformity over moving moduli,
masks and frozen labels. Here $q$ is OpenAI's bound for the moving conductor support. The zero-free
hypothesis is needed only for the finite-order slot characters, so $\kappa=3/4$ is allowed.
The changes to the proof that we can identify are four.
\begin{itemize}
\item[(a)] \emph{Hecke reflection.} This uses Lemma~\ref{lem:cm-growth}. The dual profile is
 $\mathcal MW^\sharp(s)=\mathcal MW(1-s)\,\Gamma(s+\frac12)/\Gamma(\frac32-s)$, in place of
 $\mathcal MW(1-s)\Gamma(s)/\Gamma(1-s)$. Both quotients grow like $|t|^{2\sigma-1}$ on vertical lines,
 and the conductor scale changes only at the prime above $3$. So the reflected lengths and
 seminorm bounds are unchanged.
\item[(b)] \emph{Poisson summation in the row.} The two finite Poisson transforms act on the row
 variable $k$ with radial weights, so $\alpha(l)$ is an inert coefficient.
\item[(c)] \emph{Main terms.} The induction of~\cite{OAI78} meets volume-sized main terms on rows whose
 character lies in $\Theta_0$, and cancels them with an equal-scale centered difference. In the CM case
 the corresponding lattice sums are $\sum_z\alpha(z)^ef(\N z)$ over a fixed lattice. By Poisson
 summation their main term is $\int_{\mathbb C}\alpha(z)^ef(|z|^2)\,dz=0$. This removes the volume main
 term. It does not by itself give the uniform error bounds or the induction.
\item[(d)] \emph{The zero-free hypothesis.} The slot amplitudes have finite order, so the hypothesis
 is OpenAI's theorem, with $\kappa=3/4$.
\end{itemize}
These changes are what we expect. They are not a proof of the CM statement. That requires redoing the
induction of~\cite{OAI78} with the coefficient $\alpha(l_1l_2)^e\,\tau_1\bigl(\prod l_1l_2\bigr)$ carried through
every transform, and with the angular factor kept off the slot primes. This is item~5 of
Section~\ref{sec:78status}.

\section{The final inequality}\label{sec:78final}
A bin of rows of norm $U=Z^d$ contributes at most $Z^{C(7/8)+E(d)+o(1)}$, with the exponent
$E(d)$ of~\cite[eq.~common-high-exponent]{OAI78} \cite[\S6]{G78}. Its terms come from four sources:
\begin{itemize}
\item the scales and the Mellin weight on the central contour;
\item the number of bad rows, given by the moments above with the same exponents;
\item the prime amplitudes, which have finite order;
\item the size $U^{\delta/2}$ of the finite-order numerator on the reflected line.
\end{itemize}
None of these involves the gamma factor of a CM $L$-function. That factor affects only fixed powers
of the height.

So the endpoint certificate of~\cite[Lemma~lem:balanced-endpoint]{OAI78},
\[
 -E(h)\ge\frac{49}{440640}\qquad(0\le\delta\le5/6,\ 0\le x\le1/2),
\]
with $q=x\delta$ and $R=R_*(\delta,x)$,
applies verbatim. \cite{OAI78} compares it with $\Delta=\beta_*-7/8$ in
\cite[eq.~adaptive-endpoint]{OAI78}:
\[
 E_{\rm actual}(h)-\Delta\le-\frac{49}{440640}-\Bigl(1-\frac h4\Bigr)\Delta+\epsilon_{\rm total}<0 .
\]
The term $h\Delta/4$ is a capacity loss caused by the choice $\kappa=2\beta_*-1$. With $\kappa=3/4$ it
disappears, which only helps.
The bin ceiling $\delta\le5/6$ comes from the $11/12$ stage. The balanced probe of
Section~\ref{sec:78probe}, at $X=Y=Z^{1/2}$ and with the sextic row count, gives
$\beta_*\le11/12$ for $\mathcal F_{\rm CM}$. This is Part~I of~\cite{OAI78} with the changes above;
it uses neither prime slots nor the plain moment. It is a statement about the whole family,
so it is not supplied by Part~I of these notes, which treats one curve at a time. Every margin is chosen before the target, as
in~\cite{OAI78}, so all margins are uniform in $\psi$. The continuation criterion then gives the
contradiction. Hence $\beta_*\le7/8$, which is Theorem~\ref{thm:78}.

\section{What has been checked, and what remains}\label{sec:78status}
We read~\cite{OAI78} in eight consecutive parts. For each part we listed every use of a property of
the target or of a row character, and classified it as unchanged, routine, or needing a new argument.
None needed a new argument. The table summarizes.
\begin{center}
\small
\begin{tabular}{p{2.7cm}p{4.6cm}p{7.2cm}}
\toprule
\cite{OAI78} lines & content & CM change\\\midrule
1--1648 & criterion, Gauss sums, growth lemmas & growth with $\Gamma(s+\frac12)$ (Lemma~\ref{lem:cm-growth});
 OpenAI's theorem for finite-order rows\\
1648--3313 & completed reflection, row energy & $\partial_{\bar z}^{\,2}\theta$; table in Section~\ref{sec:78probe}\\
3313--4687 & probe, Poisson rows, detector & weight~\eqref{eq:cm-weight}; local identity formal in $\psi(p)$;
 label set $\mathcal X_u$\\
4687--6807 & sextic count, high side, $11/12$ & none beyond the above; principal-denominator cases disappear\\
6807--9209 & compensated probe, low side & normalization $\overline{\psi(p)}$; slots unchanged\\
9209--12477 & inverse moment & co-oriented classes $\mu\overline\alpha\leftrightarrow\overline\alpha^{\,2}\gamma_2$ (and conjugates)\\
12477--14985 & plain fourth moment & reflection with $\Gamma(s+\frac12)$; lattice main terms vanish\\
14985--16677 & row counts, endpoint & none; $\kappa=3/4$\\\bottomrule
\end{tabular}
\end{center}

Six statements must still be written out in full, at the generality of~\cite{OAI78}:
\begin{enumerate}
\item the twisted transformation formula for $\partial_{\bar z}^{\,2}\theta$ (Appendix~\ref{app:reflection}),
 for a general modulus supported on $S$, general log-Schwartz tests, and the cross-prime phase
 bookkeeping of~\cite[L2003--2279]{OAI78}; together with its marked form
 (\cite[Cor.~cor:marked-reflection]{OAI78}) and the lemmas of~\cite{OAI78} built on it for the
 weight $\overline\alpha^{\,2}$: lem:hybrid-energy, lem:reflected-energy and lem:probe-row-norm;
\item Lemma~\ref{lem:cm-growth} together with the CM form of the logarithmic control;
\item the detector with the label set $\mathcal X_u$;
\item the inverse moment~\eqref{eq:cm-inverse} for co-oriented coefficients;
\item the plain fourth moment with $\alpha^e$ on the plain variables;
\item the $11/12$ stage of~\cite{OAI78} (its Part~I) for $\mathcal F_{\rm CM}$, giving $\beta_*\le11/12$.
\end{enumerate}
Each is obtained from the corresponding statement of~\cite{OAI78} by the changes recorded above. We
checked the new identities numerically:
\begin{itemize}
\item the point counts against $\psi_d^*$;
\item \eqref{eq:signal-cm} and its companion;
\item the vanishing of $\int\alpha^ef$;
\item the Bessel--Mellin formulas;
\item the reflection constants;
\item the Stirling exponents of all the gamma quotients.
\end{itemize}
Until the six statements are written out and refereed, Theorem~\ref{thm:78} should be read as a
detailed plan of proof, not as an established result.

\section{CM newforms of higher weight: why the method should extend}\label{sec:higherweight}

Write $\psi_k=\overline\alpha^{\,k}\eta$ with $k\ge1$ and $\eta$ of finite order. These are the unitary Hecke
characters of $K$ with infinity type $(\bar z/|z|)^k$. Up to finitely many Euler factors, their
$L$-functions are those of the CM newforms of weight $k+1$ with complex multiplication by $K$, shifted
to the center $1/2$: if $f$ is such a form, then $L(f,s+\frac k2)=L(s,\psi)$ for $\psi=\psi_k$ or its
conjugate. This section explains why we expect the following.

\paragraph{Expected statement.} For every $k\ge1$ and every finite-order $\eta$, $L(s,\overline\alpha^{\,k}\eta)\ne0$ for
$\Re s>7/8$. Equivalently, every CM newform $f$ of weight $k+1$ with CM by $K$ satisfies
\[
 L(f,s)\ne0\qquad\Bigl(\Re s>\frac k2+\frac78\Bigr),
\]
a zero-free strip along the right edge of the critical strip $\frac k2<\Re s<\frac k2+1$. The $11/12$
method of Part~I should give the weaker boundary $\frac k2+\frac{11}{12}$ in the same way.

We do not claim this as a theorem. The rest of this section records the changes in the number
$j=k+1$ of theta derivatives and in the Hecke Gamma factor. For each fixed $j$, the formulas have
the same structure and give the same length exponents; constants may depend on $j$.
The uniform analytic estimates listed at the end still have to be proved for each fixed weight.

\subsection{Two examples}
We computed the $q$-expansions of two eta products to $q^{3000}$ and compared them with Hecke characters.
In the exact $L$-function identities below, a product of Hecke characters means its primitive
inducing character, with all its Euler factors restored. Write $L^{(6)}$ when the Euler factors
at primes above $2$ and $3$ are deleted, as in the earlier sums.
\begin{itemize}
\item $f=\eta(3z)^8$, of weight $4$ and level $9$ ($k=3$). For all $207$ split rational primes
 $p=\pi\bar\pi<3000$, with both primary generators $\pi,\bar\pi$ tested ($414$ checks),
 $a_p=2\Re(\pi^3)$. For $p\equiv2\pmod3$, $a_p=0$ and $a_{p^2}=-p^3$. Thus
 $L(f,s+\frac32)=L(s,\alpha^3)$, and the target is the conjugate $\overline\alpha^{\,3}$, with
 $\eta$ trivial. This is the CM analogue of the case $u=1$ of the zeta argument. The expected strip
 is $19/8<\Re s<5/2$ (center $2$).
\item $f=\eta(2z)^3\eta(6z)^3$, of weight $3$ and level $12$ ($k=2$). For all $207$ split rational primes below
 $3000$, again testing both primary generators ($414$ checks),
 $a_p=2\Re\bigl(\pi^2(2/\pi)_3\bigr)=2\Re\bigl(\pi^2\chi_\pi(4)\bigr)$. For $p\nmid12$ with $p\equiv2\pmod3$,
 $a_p=0$ and $a_{p^2}=p^2$, as required by the nebentypus $\chi_{-3}$. The exclusion matters:
 $a_4=0$, not $4$. We have $L(f,s+1)=L(s,\alpha^2\chi_\bullet(4))$, and
 the target is $\overline\alpha^{\,2}\,\overline{\chi_\bullet(4)}$. The expected strip is $15/8<\Re s<2$ (center $3/2$).
 Since $a_3=-3$, the normalized Euler factor at $3$ is $(1+3^{-s})^{-1}$; the factor at $2$ is $1$.
 Thus, in the convention of the earlier sums,
 \[
  L(f,s+1)=(1+3^{-s})^{-1}L^{(6)}(s,\alpha^2\chi_\bullet(4)).
 \]
\end{itemize}

\subsection{What changes with \texorpdfstring{$k$}{k}}
\paragraph{1. The weight and the rows.} With target $\psi=\overline\alpha^{\,k}\eta$ the weight~\eqref{eq:cm-weight} becomes
\[
 \gamma_2(c)\overline{\alpha(cn^3)}\,\psi(cn^3)=\gamma_2(c)\,\overline{\alpha(cn^3)}^{\,k+1}\eta(cn^3).
\]
Proposition~\ref{prop:cm-rows} is formal in $\psi(p)$, so the rows are again
$\zeta_K^S(6z)L^S(w,\chi_\bullet(u))/L^S(x,\psi\overline{\chi_\bullet(u)})$. The numerators still have finite order,
the denominators have infinity type $(\bar z/|z|)^k$, and the slot amplitudes are unchanged.

\paragraph{2. The theta step with $j=k+1$ derivatives.} Every formula of Section~\ref{sec:theta} and of the
table in Section~\ref{sec:78probe} has a version for all $j\ge1$, as follows.
\begin{itemize}
\item \emph{Chain rule.} At $z=0$ the coordinate $z'$ is affine in $\bar z$, while $\bar z'$ and $v'$ do not
 depend on $\bar z$. Hence, exactly,
 \[
  \left.\partial_{\bar z}^{\,j}F(z',\bar z',v')\right|_{z=\bar z=0}
  =\bigl(-(cv)^{-2}\bigr)^{j}\,(\partial_{z'}^{\,j}F)\bigl(-\delta/c,\,1/(\N(c)v)\bigr).
 \]
 Every derivative of order $j\ge1$ kills the constant term at the identity cusp. We checked this
 identity symbolically for $j\le4$.
\item \emph{Direct Mellin transform.} For a nonzero Fourier mode and $\Re s>\frac16-\frac j2$,
 \[
  \int_0^\infty(2\pi i\bar\ell)^j\,vK_{1/3}(4\pi|\ell|v)\,v^{2s+j-2}\,dv
  =\frac{i^j}4(2\pi)^{-2s}\,\Gamma\Bigl(s+\frac j2-\frac16\Bigr)\Gamma\Bigl(s+\frac j2+\frac16\Bigr)
   \overline{\alpha(\ell)}^{\,j}\N(\ell)^{-s}.
 \]
 The integral converges absolutely in this half-plane; elsewhere the right side gives its
 meromorphic continuation. For $j=1,2$ this gives the two entries of the table. We checked it numerically for $j\le5$ (relative
 error below $10^{-25}$).
\item \emph{Reflected side.} The same substitution gives $\prod_\pm\Gamma(\frac{j+2}2-s\pm\frac16)$, the cusp
 factor $\overline{\alpha(c)}^{\,2j}$, and the dual angular factor $\alpha(\mu)^j$.
\item \emph{Kernel.} Write
 \[
  R_j(t)=\prod_\pm\frac{\Gamma(\frac{j+1}2+t\pm\frac16)}{\Gamma(\frac{j+1}2-t\pm\frac16)} .
 \]
 Then $R_1$ and $R_2$ are the two kernels of the table. $R_j$ is holomorphic for
 $\Re t>\frac16-\frac{j+1}2$, which is to the left of $-1/4$ for every $j\ge1$, and
 $|R_j(\sigma+iu)|\ll_{j,A}(1+|u|)^{4\sigma}$ on $-1/4\le\sigma\le A$ (checked numerically for $j\le5$). So the
 reflected weight obeys the same bound $\ll\min(x^{1/4},x^{-A})$, with constants depending on $j$.
 There is no uniform bound as $j$ grows: for example,
 \[
  R_j(1)=\frac{(3j-4)(3j-2)(3j+2)(3j+4)}{1296}\sim\frac{j^4}{16}.
 \]
\end{itemize}
The dual lengths are unchanged, because the reflected argument $\N(\mu)X/\N(c)^2$ does not depend on
$j$. In the quadratic large sieve, $\alpha(\mu)^j$ is a unimodular coefficient of the dual index. Hence
Proposition~\ref{prop:cm-completed} and the low estimates $Z^{1/4+\varepsilon}$ and $Z^{3/16+\varepsilon}$ should hold
for every $k$.

\paragraph{3. Growth.} The functional equation has $\Gamma(s+\frac k2)$ in place of $\Gamma(s+\frac12)$. The
exponent of $|\Gamma(1-s+\frac k2)/\Gamma(s+\frac k2)|$ on $\Re s=\sigma$ is $1-2\sigma$, independent of $k$, so
Lemma~\ref{lem:cm-growth} keeps its exponents. The theta series now has the harmonic polynomial
$\bar a^{\,k}$. Its Poisson transform is a multiple of $\bar\xi^{\,k}e^{-c'|\xi|^2/v}v^{-1-k}$, which vanishes at $\xi=0$,
so $L(s,\psi)$ is entire.

\paragraph{4. Moments.} In the inverse moment the co-oriented classes become
$\mu\overline\alpha^{\,k}\leftrightarrow\overline\alpha^{\,k+1}\gamma_2$. This follows from the same identity
$\mu\gamma_{-1}=\chi(-1)G^{-1}\overline\alpha\gamma_2$, and the cube factor becomes $\overline\alpha^{\,3(k+1)}$. The theta steps
use $\partial_{\bar z}^{\,k+1}\theta$. In the plain moment, Hecke reflection uses
$\Gamma(s+\frac k2)/\Gamma(1-s+\frac k2)$, a Hankel kernel of order $k$ with the same vertical growth. The
lattice main terms vanish because $\int\alpha(z)^{\pm k}f(|z|^2)\,dz=0$.

\paragraph{5. Unchanged.} The following carry over verbatim:
\begin{itemize}
\item the label sets, with $\overline\alpha^{\,k}\Theta_0\overline{\chi_\bullet(u)}$ and $\alpha^k\Theta_0\chi_\bullet(u)$;
\item the slot normalization $\overline{\psi(p)}$;
\item $\kappa=3/4$ from OpenAI's theorem;
\item every exponent in the final inequality, and hence the endpoint certificate.
\end{itemize}
For Part~I the targets are rows of the same kind. In the orientation $\muk(n)\alpha(n)^k\chi_n(u_0)$ of
Part~I, the row is $u_0=1$ for $\eta(3z)^8$ and $u_0=4$ for $\eta(2z)^3\eta(6z)^3$.

\subsection{Remarks}
\begin{itemize}
\item \emph{$k=0$ is OpenAI's case.} There $j=1$ and the weight is $\gamma_2\overline\alpha\eta$. The recipe ``target
 $\overline\alpha^{\,k}\eta$, $k+1$ theta derivatives'' therefore covers unitary Hecke characters of $K$ with
 purely angular infinity type; negative $k$ reduces to positive $k$ by conjugation. A general unitary
 Hecke character may also have a continuous norm factor $\N^{it}$, with $t\in\mathbb R$. Since
 \[
  L(s,\psi\N^{it})=L(s-it,\psi),
 \]
 these vertical shifts leave the proposed zero-free half-plane unchanged.
\item \emph{Orientation.} The conjugate target $\alpha^k\bar\eta$ would leave $\alpha^{k-1}$ in the weight. For $k=1$
 that is no angular factor at all: the theta function without a derivative, whose constant terms are
 not controlled. That is why Part~II uses $\overline\alpha\eta$. For $k\ge2$ the factor $\alpha^{k-1}$ is produced by
 $\partial_z^{\,k-1}$. The chain rule is symmetric: at $z=\bar z=0$, $\partial_z^{\,j}$ becomes
 $(-(\bar cv)^{-2})^j\partial_{\bar z'}^{\,j}$, exactly; we checked this symbolically for $j\le3$. So for $k\ge2$
 either orientation should work; we use $\overline\alpha^{\,k}\eta$ for uniformity.
\item \emph{Symmetric powers.} For a CM curve $E_d$ and $m\ge1$, put
 $L_{\rm an}(\mathrm{Sym}^mE_d,s)=L(\mathrm{Sym}^mE_d,s+m/2)$.
 Up to finitely many Euler factors, the factorization is
 \[
  L_{\rm an}(\mathrm{Sym}^mE_d,s)
  =\prod_{0\le i<m/2}L(s,(\psi_d^*)^{m-2i})
   \begin{cases}
    1,&m\text{ odd},\\
    L(s,\chi_{-3}^{\,m/2}),&m\text{ even}.
   \end{cases}
 \]
 The last factor is $\zeta(s)$ when $m\equiv0\pmod4$ and $L(s,\chi_{-3})$ when
 $m\equiv2\pmod4$, with the trivial character interpreted primitively. This is the CM
 factorization of~\cite[\S3.5]{MW}. The positive-power factors are of the type above; the remaining
 factor is covered by OpenAI's finite-order result. Together these would give a zero-free
 half-plane $\Re s>7/8$ in analytic normalization. When $4\mid m$, the function has a pole at $s=1$:
 ``zero-free'' here is a statement about a meromorphic function, not a claim of holomorphy throughout
 the half-plane. These $L$-functions are not primitive for $m\ge2$.
\item \emph{What would still have to be written.} The six restatements of Section~\ref{sec:78status} are
 needed for each $k$, with $\partial_{\bar z}^{\,k+1}$ in place of $\partial_{\bar z}^{\,2}$. We see no step where $k$
 enters other than through $j=k+1$ and the Gamma factors. The constants depend on $k$, but $k$ is fixed.
\end{itemize}
The numerical checks of this section were made with short Python scripts.

\appendix
\section{The reflected series and explicit theta coefficients}\label{app:reflection}

This appendix supplies the coefficient data for readers who want every
series in the theta step in finite-sum terms. It can be skipped on a
first reading. The imported theta symmetry and coefficient tables are
those of Kubota and Patterson recorded in~\cite[Section 5.1 and
Appendix]{DR}, in the conventions of~\cite{OAI78}.
Only these theta facts are used from~\cite{DR}; its GRH-conditional
prime Gauss-sum asymptotic is not an input here.

For $\Psi(n)=\xi_0(n)\prod_{\pi\in\mathcal P}\chi_\pi(n)^{j_\pi}$,
the transformation expresses $T(X;\Psi)$ as a finite sum of terms
\begin{equation}\label{eq:reflection}
 C\sum_{0\ne\ell\in\lambda^{-4}\OO}
 \frac{d_\sigma(\ell)\alpha(\ell)^2}{\sqrt{\N(\ell)}}
 \psi(\lambda^4\ell)\prod_{\pi\in\mathcal A}B_{\pi,j_\pi}(\lambda^4\ell)
 V^\sharp\left(\frac{\N(\ell)X}{\N(c)^2}\right).
\end{equation}
Here $\sigma\in\{0,+,-\}$,
$\{\pi:j_\pi\ne0\}\subseteq\mathcal A\subseteq\mathcal P$,
$c=c_0\prod_{\pi\in\mathcal A}\pi$,
$c_0\mid8\lambda^2$ up to a unit, and
$\psi(x)=e(\beta x/(\lambda^3c_0))$ for an Eisenstein integer $\beta$.
There are at most $5184\cdot2^{|\mathcal P|}$ terms, each with
$|C|\le1/81$. The constants are obtained by finite additive Fourier
expansion of $\Psi$ and the theta transformation, as in the completed reflection
of~\cite[Prop.~prop:completed-reflection]{OAI78} with two derivatives in place of one.
The expression is an absolutely convergent series. The weights
$V^\sharp$ and the local factors $B_{\pi,j}$ were explicitly defined
in~\eqref{eq:weight} and~\eqref{eq:local}.

For the remaining coefficients, define the finite cubic Gauss sum
\[
 g(r,c)=\sum_{x\bmod c}\left(\frac xc\right)_3\ce(rx/c).
\]
In this appendix, $c,b$ are primary, $c$ is squarefree, and they may be
even. The cubic symbol at the even primary prime $-2$ is defined by
$(x/(-2))_3\equiv x\pmod2$ for $x$ prime to $2$, and zero otherwise;
at other primary primes use the exponent $(\N(\pi)-1)/3$.
Put $\zeta_9=e^{2\pi i/9}$. Define $\tau(\ell)$ by the following table
and set it equal to zero at every other index:
\[
\begin{array}{c|c}
 \ell&\tau(\ell)\\\hline
 \pm\lambda^{3\nu-4}cb^3,\ \nu\ge1
 &3^{\nu/2+2}\overline{g(\lambda^2,c)}\,|b/c|\\
 \pm\omega\lambda^{3\nu-4}cb^3,\ \nu\ge1
 &3^{\nu/2+2}\zeta_9^{-1}\overline{g(\omega\lambda^2,c)}\,|b/c|\\
 \pm\omega^2\lambda^{3\nu-4}cb^3,\ \nu\ge1
 &3^{\nu/2+2}\zeta_9\overline{g(\omega^2\lambda^2,c)}\,|b/c|\\
 \pm\lambda^{3\nu-3}cb^3,\ \nu\ge0
 &3^{\nu/2+5/2}\overline{g(1,c)}\,|b/c|
\end{array}
\]
For the two other cusps, define $\tau_1,\tau_2$ by
\[
\begin{array}{c|c|c}
 j&\tau_1(\omega^j\lambda^{-4}cb^3)
   &\tau_2(-\omega^j\lambda^{-4}cb^3)\\\hline
0&9\omega\overline{g(\lambda^2,c)}|b/c|
 &9\omega^2\overline{g(\lambda^2,c)}|b/c|\\
1&9\zeta_9^{-1}\omega^2\overline{g(\omega\lambda^2,c)}|b/c|
 &9\zeta_9^{-1}\overline{g(\omega\lambda^2,c)}|b/c|\\
2&9\zeta_9\overline{g(\omega^2\lambda^2,c)}|b/c|
 &9\omega\zeta_9\overline{g(\omega^2\lambda^2,c)}|b/c|
\end{array}
\]
and zero elsewhere. Finally set
\[
 t_0(\ell)=\tau(\ell),\qquad
 t_+(\ell)=\omega\tau_2(\omega\ell)\ce(\ell),\qquad
 t_-(\ell)=\omega^2\tau_1(\omega^2\ell)\ce(\ell),\qquad
 d_\sigma(\ell)=\overline{t_\sigma(-\ell)}.
\]
These are precisely the coefficients of the conjugated theta functions
at the three cusps. Their support is contained in the set of indices
$\ell=u\lambda^mnb^3$, with $u$ a unit, $m\ge-4$, $n$ squarefree,
and $n,b$ primary (the tables above impose further restrictions), and
\[
 |d_\sigma(u\lambda^mnb^3)|\le27\cdot3^{m/6}|b|.
\]
In the varying quadratic symbol in~\eqref{eq:reflection}, the factor
$b^3$ contributes $(b/k)_2$, since $(b^3/k)_2=(b/k)_2$.
Fixing $u,m,b$ leaves squarefree columns $n$, exactly as required for
\eqref{eq:qls}. This is why the explicit support of the theta
coefficients matters to the averaging argument.


\begin{thebibliography}{9}
\bibitem{ZG} \emph{Sixth powers in a mean square: an accessible guide to
 the $11/12$ argument for $\zeta(s)$}, companion notes,
 \url{https://www.math.uwaterloo.ca/~mrubinst/quasi_rh_notes/zeta_11_12_guide.pdf}.
\bibitem{OAI1112} OpenAI, \emph{The Quasi-Riemann Hypothesis}, manuscript
 dated October 5, 2026.
 \href{https://github.com/openai/math/tree/main/preprints/The-Quasi-Riemann-Hypothesis-October-5-2026}{Manuscript and source on GitHub}.
\bibitem{OAI78} OpenAI, \emph{The Quasi-Riemann Hypothesis: A Zero-Free
 Half-Plane $\Re s>7/8$}, manuscript dated September 30, 2026,
 especially the main theorem and the definition of $\beta_*$.
 Labels and line numbers refer to its source file \texttt{build/paper.tex}.
 \href{https://github.com/openai/math/tree/main/preprints/The-Quasi-Riemann-Hypothesis-September-30-2026}{Manuscript and source on GitHub}.
\bibitem{Neukirch} J. Neukirch, \emph{Algebraic Number Theory},
 Grundlehren der mathematischen Wissenschaften 322, Springer, 1999.
\bibitem{G78} \emph{Zeros of a whole family: an accessible guide to the
 $7/8$ argument for $\zeta(s)$}, companion notes,
 \url{https://www.math.uwaterloo.ca/~mrubinst/quasi_rh_notes/zeta_7_8_guide.pdf}.
\bibitem{IR} K. Ireland and M. Rosen, \emph{A Classical Introduction to Modern Number Theory},
 2nd ed., Graduate Texts in Mathematics 84, Springer, 1990.
\bibitem{MW} P. Martin and M. Watkins, \emph{Symmetric powers of elliptic curve $L$-functions},
 2006, \S3.5.
 \url{https://arxiv.org/abs/math/0604095}.
\bibitem{DR} A. Dunn and M. Radziwi\l\l,
 \emph{Bias in cubic Gauss sums: Patterson's conjecture},
 Annals of Mathematics (2) \textbf{200} (2024), 967--1057.
 \url{https://arxiv.org/abs/2109.07463}.
\bibitem{GL} L. Goldmakher and B. Louvel,
 \emph{A quadratic large sieve inequality over number fields},
 Mathematical Proceedings of the Cambridge Philosophical Society
 \textbf{154} (2013), 193--212.
 \url{https://arxiv.org/abs/1112.1642}.
\end{thebibliography}
\end{document}
